\chapter{Smooth varieties in arbitrary characteristic}
\index{positive characteristic!envelope theorem}\index{test ideal}
\index{multiplier ideal!asymptotic}\index{envelope property!smooth case}
\label{chap:bj-smooth-characteristic}

\section{Orientation}
\label{sec:ch40-orientation}

\subsection{What is proved here, and against what}
\label{sec:ch40-what}

Let $X$ be a smooth projective variety over the algebraically closed
trivially valued field $k$, of arbitrary characteristic, and let
$\theta\in N^1(X)_{\mathbf R}$.  This chapter proves that
$\theta$ has the envelope property of
\Cref{def:bj-envelope-property}:
\[
 \begin{gathered}
 (\phi_\alpha)_{\alpha\in A}\subset\PSH_X(\theta)
 \ \text{nonempty and bounded above}\\
 \Longrightarrow\quad
 \Bigl(\sup_\alpha\phi_\alpha\Bigr)^*\in\PSH_X(\theta).
 \end{gathered}
\]
This is \Cref{thm:bj-smooth-envelope}.  It is the only step of
\Crefrange{chap:bj-cartier-realization}{chap:bj-envelope-consequences} at
which a genuinely new psh function has to be \emph{manufactured}: every
other chapter of the proof transports psh functions along a morphism, and
transport can only move a function that already exists.  Here there is no
morphism and nothing to transport.  The function has to be built out of
sections of line bundles, and the mechanism that builds it is
ideal-theoretic.

The mechanism is easy to state and its shape should be held in view
throughout.  One is given an obstacle $h$, a continuous function on
$X^{\An}$ below which the psh competitors are constrained to lie, and one
must show that the largest competitor is itself psh.  The competitors are
suprema, and \Cref{thm:bj-basic-closure} says nothing about suprema.  The
device is to replace the increasing family of competitors by a
\emph{decreasing} family of explicitly constructed Fubini--Study functions
with the same limit.  Decreasing families are what
\Cref{def:bj-psh} is built from, so the limit is psh for free once the
family has been constructed and once one knows that the limit is not
identically $-\infty$.

Producing that decreasing family is the whole difficulty, and it is done in
three moves.

\emph{First}, the obstacle is turned into a graded sequence of ideals.
Writing the obstacle as $\phi_D$ for a vertical divisor $D\le0$ on a
model $\mu:\cY\to X\times\mathbf A^1$ of the trivial test configuration,
one forms the relative base ideals $\mathfrak a_m$ of $m(\mu^*L+D)$ and
pushes them to $X\times\mathbf A^1$.  The resulting ideals $J_m$ satisfy
$J_mJ_r\subseteq J_{m+r}$, and the associated normalized ideal functions
$m^{-1}\phi_{J_m}$ increase to the largest Fubini--Study function below
the obstacle.  Increasing: this is still the wrong direction.

\emph{Second}, each $J_m$ is replaced by a larger ideal
$B_m\supseteq J_m$ chosen so that the assignment $m\mapsto B_m$ is
\emph{sub}additive:
$B_{m\ell}\subseteq B_m^{\,\ell}$.  Under the order-reversing dictionary
$I\mapsto\log|I|=-v(I)$ of \eqref{eq:ch34-log-ideal}, subadditivity of
ideals becomes monotonicity of functions, and the normalized functions
$u_m:=m^{-1}\phi_{B_m}$ \emph{decrease} along the divisibility order.  In
characteristic zero, $B_m$ is the asymptotic multiplier ideal
$\mathcal J_S(J_\bullet^m)$; in characteristic $p>0$ it is the
asymptotic test ideal $\tau_S(J_\bullet^m)$.  These are different sheaves
constructed by different means, and the chapter does not pretend otherwise.

\emph{Third}, one must check that the decreasing family
$(u_m)$ consists of Fubini--Study functions in classes converging to
$c_1(L)$, and that its limit has the right values.  The first is a global
generation statement with a twist $A$ that does \emph{not} depend on
$m$, so that the stage classes are exactly $c_1(L)+m^{-1}c_1(A)$; the
second is a valuation estimate with an error $A_S(w)/m$ at each fixed
divisorial valuation $w$.  Both errors are $O(m^{-1})$, both are
tracked explicitly below, and they are \emph{different errors}: one lives in
$N^1(X)$ and one lives in $\mathbf R$.  Conflating them is the standard
way to produce a proof that looks complete and is not.

\subsection{One architecture, two engines}
\label{sec:ch40-two-engines}

The point of the chapter, and the reason it is written at this length, is
that the three moves above run \emph{verbatim} in every characteristic once
the ideal engine has been fixed.  What changes is the engine, and the honest
description of the change is not ``multiplier ideals become test ideals'' but
the following table, whose four rows are the only ideal-theoretic inputs the
argument consumes.

\begin{center}
\begin{tabular}{p{0.20\textwidth}p{0.34\textwidth}p{0.34\textwidth}}
\toprule
required output & $\operatorname{char}k=0$ & $\operatorname{char}k=p>0$\\
\midrule
containment
 & $J_m\subseteq\mathcal J_S(J_\bullet^m)$, from the strict inequality in
   the valuative criterion
 & $J_m\subseteq\tau_S(J_\bullet^m)$, from the trace description\\
subadditivity
 & Demailly--Ein--Lazarsfeld: restriction to the diagonal of
   $S\times S$
 & Musta\c t\u a: a direct Frobenius computation; no diagonal, no
   restriction theorem\\
fixed twist
 & relative Kawamata--Viehweg vanishing on a log resolution
 & Frobenius trace, the projection formula, and Serre vanishing; no
   resolution\\
valuation error
 & $v(\mathcal J)\ge m\,v(J_\bullet)-A_S(v)$ from the valuative
   description
 & $w(\tau)\ge t\,w(I)-A_S(w)$ from Canton's Berkovich log
   discrepancy\\
\bottomrule
\end{tabular}
\end{center}

Three observations about this table are the mathematical content of
\Cref{sec:ch40-comparison}, and they are worth stating at the outset so that
the reader knows what to watch for.

\emph{The two subadditivity theorems are genuinely different theorems.}  In
characteristic zero, subadditivity is proved by realizing
$\mathcal J(\mathfrak a\mathfrak b)$ as the restriction to the diagonal of
an external product on $S\times S$, and the proof uses a restriction
theorem for multiplier ideals which is a corollary of Nadel vanishing.  In
characteristic $p$ there is no such route, and the proof is a direct
manipulation of Frobenius roots.  The two statements coincide; the two
proofs share no step.

\emph{The Frobenius input has no characteristic-zero counterpart.}  The
surjection \eqref{eq:bj-smooth-trace-quotient} from a Frobenius pushforward
onto $\tau\otimes\omega_S$ is what makes a twist independent of $m$
available in characteristic $p$; the residual coefficient
\eqref{eq:bj-smooth-residual-coefficient} lies in a fixed finite set while
the positive coefficient $p^e(d-i)$ of the ample bundle grows without
bound, and this asymmetry is the whole reason a single $A$ works.  In
characteristic zero nothing plays this role: the same conclusion is reached
through a vanishing theorem on a log resolution, and log resolutions are not
available in characteristic $p$ in any dimension for which the argument
would need them.

\emph{No comparison between the two engines is used.}  The theorem of
Hara--Yoshida identifying multiplier ideals with test ideals after reduction
modulo $p$, recorded as \Cref{thm:appC-hara-yoshida}, is a comparison for
one fixed datum after excluding finitely many primes.  It is not used here
and could not be: the graded sequence $J_\bullet$ varies with $m$, the
ground field is fixed, and there is nothing to reduce.  The two arguments
are run separately from beginning to end.  \Cref{ex:ch40-cusp} exhibits a
graded sequence for which the two engines produce genuinely different ideals,
so this separation is not a stylistic preference.

\subsection{The firewall}
\label{sec:ch40-firewall}

This chapter is one of the inputs of \Cref{thm:bj-final-envelope}: the proof
map of \Cref{sec:ch37-proof-map} lists it as step four, and
\Cref{chap:bj-alterations} reduces the normal case to it through
\Cref{chap:bj-finite-descent}.  Any appeal here to the envelope theorem, or
to any of its immediate consequences, would close a loop through
\Crefrange{chap:bj-cartier-realization}{chap:bj-envelope-consequences}.

\begin{warning}{The envelope theorem is not available in this chapter}
\label{warn:ch40-firewall}
Nothing in \Cref{chap:bj-smooth-characteristic} may use, directly or
transitively:
\begin{enumerate}
\item \Cref{thm:bj-final-envelope}, in any of its formulations, nor
      \Cref{def:bj-envelope-property-final};
\item the envelope property $\operatorname{EP}$ of
      \Cref{def:bj-envelope-property} as a \emph{theorem} for any class on
      any variety; below it is either the conclusion being proved, or --- in
      \Cref{lem:ch40-ep-nef-limit} alone --- an explicit displayed
      hypothesis about \emph{other} classes on the same $X$, discharged
      before it is used;
\item closure of $\PSH_X(\theta)$ under suprema, regularized or not;
\item the obstacle envelope $P_\theta$ of
      \Cref{def:bj-obstacle-envelope} regarded as an operator with values
      in $\PSH_X(\theta)$.  Throughout this chapter $P_L(\phi_D)$
      denotes \eqref{eq:bj-smooth-obstacle-envelope-definition}, the
      upper-semicontinuous regularization of the raw supremum
      \eqref{eq:bj-obstacle-envelope}; neither function is known to be psh
      until \eqref{eq:bj-smooth-obstacle-envelope} is proved, and
      \Cref{warn:ch40-two-envelopes} records when the two coincide;
\item \Cref{thm:ch35-cartier-psh-criterion} in its direction
      $\Leftarrow$, and its consequences
      \Cref{thm:bj-nef-bdivisor-correspondence} and
      \Cref{cor:ch35-monotone-approximation};
\item any statement of \Cref{chap:bj-envelope-consequences}.
\end{enumerate}
Every limit taken below is a decreasing limit along the divisibility order.
Every supremum that appears is either the increasing base-ideal supremum
\eqref{eq:bj-smooth-Q}, whose limit $Q$ is never claimed to be psh, or the
raw obstacle supremum of item~(4), which is compared with a psh function
only through the one-sided comparison
\Cref{prop:ch39-determination-package}\,\cref{itm:ch39-det:compare}, in which
the psh function is the \emph{smaller} one.
\end{warning}

Two clarifications keep the firewall usable.

First, \Cref{thm:general-obstacle-criterion} is not forbidden.  It is a
formal equivalence, proved in \Cref{chap:bj-envelope-history} from
divisorial nonnegligibility and the Alexander--Taylor bound, between the
envelope property for a fixed pair $(X,\theta)$ and the psh-ness of
$P_\theta(h)$ for PL obstacles $h$.  We use only the implication
``obstacle envelopes are psh $\Rightarrow$ $\operatorname{EP}$'', which
is the direction that consumes nothing.

Second, \Cref{thm:bj22-nef-envelope} and
\Cref{thm:bj-addendum-all-classes} --- the historical results of
Boucksom--Jonsson and their addendum --- are \emph{not} used as inputs,
although the present theorem contains them.  They are recorded in
\Cref{chap:bj-envelope-history} as history and are proved there under
hypotheses (characteristic zero, or dimension at most two) that this chapter
removes.  Citing them here would be harmless for the characteristic-zero
statement and useless for the general one; the argument below is written so
that a reader can check that no such citation occurs.

\begin{remark}[Why the firewall is not vacuous here]
\label{rem:ch40-firewall-not-vacuous}
The temptation in this chapter is concrete and it is not the same as the one
described in \Cref{rem:ch39-firewall-not-vacuous}.  It is this.  Once the
envelope theorem is known, $P_L(\phi_D)$ is a psh function, and the entire
ideal-theoretic construction below can be replaced by the sentence
``$P_L(\phi_D)$ is psh by \Cref{thm:bj-final-envelope}, and equals $Q$ on
divisorial points by \eqref{eq:bj-smooth-off-base}''.  A reader who has the
finished theory in view may therefore feel that the ideals are an elaborate
way of proving something obvious.  They are not: without the theorem,
$P_L(\phi_D)$ is a pointwise supremum of an uncountable family, its usc
regularization is a limsup over neighbourhoods in a non-metrizable compact
space, and there is no reason for either to be a decreasing limit of
Fubini--Study functions.  The ideals produce, by hand, the decreasing family
that the definition of $\PSH$ requires and that the supremum does not
supply.
\end{remark}

\subsection{The six tasks, and where each is discharged}
\label{sec:ch40-tasks}

\begin{enumerate}
\item\label{itm:ch40-task:reduction}
 \emph{Reduce to a rational obstacle in a big rational class.}
 \Cref{sec:ch40-reduction}: real classes are approximated from above by
 rational big ones with nef differences, arbitrary families are replaced by
 PL obstacles, and PL obstacles are realized on models of the trivial test
 configuration.  The conclusion is
 \Cref{prop:bj-smooth-obstacle-interface}.
\item\label{itm:ch40-task:ideals}
 \emph{Convert the obstacle into a graded sequence of ideals.}
 \Cref{sec:ch40-base-ideals}: the relative base ideals of $m\cL$, their
 pushforward $J_m$ to $X\times\mathbf A^1$, the base-ideal identity
 \eqref{eq:bj-smooth-obstacle-identity}, the two gradings, and the complete
 base-ideal limit $Q$.  \Cref{ex:ch40-p2-obstacle} computes all of this
 for an explicit obstacle on $\mathbf P^2$.
\item\label{itm:ch40-task:engines}
 \emph{Install the two ideal engines.}  \Cref{sec:ch40-engines}: the four
 properties of the table above, with multiplier ideals in characteristic
 zero and asymptotic test ideals in characteristic $p$; the appendix
 \Cref{app:ideals} supplies the constructions, this section supplies the
 deductions actually used.
\item\label{itm:ch40-task:twist}
 \emph{Fixed twist and exact class error.}  \Cref{sec:ch40-stages}: the
 twist $A=K_X+dH$, the relative generation statements in both
 characteristics, the identification of the stage class as exactly
 $c_1(L)+m^{-1}c_1(A)$, and the passage to the limit.
\item\label{itm:ch40-task:valuation}
 \emph{The valuation inequality and the psh limit.}
 \Cref{sec:ch40-valuation}: \Cref{lem:bj-smooth-canton} and its
 characteristic-zero counterpart, the divisorial identification
 \Cref{prop:bj-smooth-divisorial-identification}, and the proof that the
 limit is $\theta$-psh rather than $\equiv-\infty$.
\item\label{itm:ch40-task:comparison}
 \emph{Compare the two proofs.}  \Cref{sec:ch40-comparison}: what is shared,
 what is not, and which steps would break if one tried to transplant an
 argument from one characteristic to the other.
\end{enumerate}

\subsection{Standing hypotheses}
\label{sec:ch40-hypotheses}

Throughout the chapter,
\begin{equation}
\label{eq:ch40-standing}
 \begin{gathered}
 k=\bar k\ \text{trivially valued},\qquad
 X\ \text{smooth projective over }k,\qquad n:=\dim X,\\
 S:=X\times\mathbf A^1,\qquad
 \rho:=\operatorname{pr}_{\mathbf A^1}:S\to\mathbf A^1,\qquad
 \varpi:=\text{the coordinate on }\mathbf A^1 .
 \end{gathered}
\end{equation}
We do \emph{not} assume $X$ connected in
\eqref{eq:ch40-standing}; but a variety is integral by the convention sheet,
so $X$ is in fact irreducible, and every ideal considered below is nonzero
in the sense of \Cref{def:appC-graded-sequence}.  The scheme $S$ is
smooth over $k$ of dimension $n+1$, hence regular; when
$\operatorname{char}k=p>0$ it is moreover $F$-finite, because $k$ is
perfect --- being algebraically closed --- and $S$ is of finite type over
$k$.  Both facts are used without further comment.

The Gauss extension of \Cref{sec:ch32-gauss} is normalized by
\begin{equation}
\label{eq:ch40-gauss}
 \sigma(v)\Bigl(\sum_i f_i\varpi^i\Bigr)=\min_i\{v(f_i)+i\},
 \qquad\text{so}\qquad
 \sigma(v)(\varpi)=1,
\end{equation}
and for a coherent ideal $I\subseteq\cO_S$ we write, extending
\eqref{eq:ch34-log-ideal},
\begin{equation}
\label{eq:ch40-ideal-function}
 \phi_I:X^{\An}\longrightarrow[-\infty,0],
 \qquad
 \phi_I(v):=-\sigma(v)(I).
\end{equation}
This is the only sign convention in the chapter that has to be memorized;
\Cref{warn:ch40-signs} records the two unit tests that fix it.

\subsection{Dependencies}
\label{sec:ch40-dependencies}

From \Cref{chap:trivial-analytification} we use the Gauss extension
\eqref{eq:gauss-extension-triv} with its transfer properties
\Cref{thm:ch32-gauss-transfer}, the density of divisorial points
\Cref{thm:triv-divisorial-density}, and the quasi-monomial calculus
\Cref{thm:ch32-quasimonomial-invariants}.  From
\Cref{chap:bj-test-configurations} we use flag ideals
\Cref{def:bj-flag-ideal} with the weight decomposition
\Cref{lem:ch33-flag-weight-decomposition}, and the base-ideal formalism of
\Cref{sec:ch33-base-ideals}.  From \Cref{chap:bj-fs-pl-psh} we use
\Cref{def:gf-fs-functions,prop:ch34-generic-finiteness,def:bj-psh,thm:bj-basic-closure},
the flag-ideal criterion \Cref{prop:ch34-fs-flag-ideal} and the test
configuration correspondence \Cref{thm:tc-pl-correspondence}.  From
\Cref{chap:bj-homogeneity-bdivisors} we use divisorial comparison
\Cref{thm:ch35-divisorial-comparison}, packaged for this part of the book as
\Cref{prop:ch39-determination-package}.  From
\Cref{chap:bj-envelope-history} we use the obstacle calculus
\Cref{lem:obstacle-envelope-formal} and the criterion
\Cref{thm:general-obstacle-criterion}.  From
\Cref{chap:bj-cartier-realization} we use the vocabulary of the two gradings
\Cref{prop:ch38-two-gradings} and the drift discipline
\Cref{prop:ch38-drift-ledger}, and from \Cref{chap:bj-finite-descent} the
pullback lemma \Cref{lem:ch41-proper-gf-pullback} and the real-weight
corollary \Cref{cor:ch41-real-weights}.  The ideal theory is
\Cref{app:ideals}; the birational inputs are
\Cref{thm:appB-resolution} in characteristic zero and nothing at all in
characteristic $p$.

\section{The rational obstacle reduction}
\label{sec:ch40-reduction}

The envelope property is a statement about an arbitrary bounded-above family
in an arbitrary real class.  Nothing ideal-theoretic can be said about such
an object directly.  This section replaces it, in four steps, by a statement
about one explicit function attached to one explicit divisor on one explicit
model: real classes become big rational ones, families become obstacles,
obstacles become PL functions, and PL functions become vertical divisors.
Each step is reversible in the sense that the conclusion of the last step
implies the conclusion of the first; none of them uses any positivity or
characteristic hypothesis on $X$ beyond projectivity and integrality.

\subsection{From real classes to big rational classes}
\label{sec:ch40-real-to-rational}

Two elementary lemmas do the work.  The first is the only place in the
chapter where the envelope property occurs as a hypothesis rather than as a
conclusion, and it is discharged immediately: the classes it assumes it for
are rational and big, which is exactly what the rest of the chapter treats.

\begin{lemma}
\label{lem:ch40-ep-nef-limit}
Let $X$ be an integral projective $k$-variety, let
$\theta\in N^1(X)_{\mathbf R}$, and let $(\theta_j)_{j\in J}$ be a family
in $N^1(X)_{\mathbf R}$ such that
\[
 \theta_j-\theta\in\Nef(X)\ \text{for every }j,
 \qquad
 \theta_j\longrightarrow\theta\ \text{in }N^1(X)_{\mathbf R}.
\]
If $\operatorname{EP}_X(\theta_j)$ holds for every $j$, then
$\operatorname{EP}_X(\theta)$ holds.
\end{lemma}

\begin{proof}
Let $(\phi_\alpha)_{\alpha\in A}\subset\PSH_X(\theta)$ be nonempty and
bounded above, and put $\phi:=(\sup_\alpha\phi_\alpha)^*$.

Fix $j$.  Since $\theta_j-\theta$ is nef,
\Cref{lem:ch34-psh-nonempty} gives $0\in\PSH_X(\theta_j-\theta)$, and
\Cref{thm:bj-basic-closure}\,(3) applied with $s=t=1$ to the pair
$(\phi_\alpha,0)$ gives
\[
 \phi_\alpha=\phi_\alpha+0
 \in\PSH_X\bigl(\theta+(\theta_j-\theta)\bigr)=\PSH_X(\theta_j).
\]
No forbidden sum occurs: the second summand is the constant $0$.  Thus
$(\phi_\alpha)_\alpha$ is a nonempty bounded-above family in
$\PSH_X(\theta_j)$, and its regularized supremum is the same function
$\phi$.  By $\operatorname{EP}_X(\theta_j)$ we get
$\phi\in\PSH_X(\theta_j)$.

Since $j$ was arbitrary, the set
$\{\eta\in N^1(X)_{\mathbf R}:\phi\in\PSH_X(\eta)\}$ contains every
$\theta_j$.  By the closedness assertion in
\Cref{thm:bj-basic-closure}\,(5) that set is closed, and $\theta_j\to\theta$;
hence $\phi\in\PSH_X(\theta)$.
\end{proof}

\begin{lemma}
\label{lem:ch40-rational-approximation}
Let $X$ be an integral projective $k$-variety and
$\theta\in\operatorname{Psef}(X)$.  Fix a norm on $N^1(X)_{\mathbf R}$.
Then for every $\delta>0$ there is a $\mathbf Q$-line bundle $L$ on
$X$ with
\[
 \lVert c_1(L)-\theta\rVert<\delta,
 \qquad
 c_1(L)-\theta\ \text{ample},
 \qquad
 c_1(L)\ \text{big}.
\]
\end{lemma}

\begin{proof}
Fix an ample $\mathbf Q$-class $\alpha$, which exists because $X$ is
projective.  The set
\[
 U:=\bigl\{\eta\in N^1(X)_{\mathbf R}\ :\
 \eta-\theta\in\Amp(X),\ \lVert\eta-\theta\rVert<\delta\bigr\}
\]
is open, because $\Amp(X)$ is an open cone in $N^1(X)_{\mathbf R}$, and
nonempty, because it contains $\theta+\varepsilon\alpha$ for
$0<\varepsilon<\delta/\lVert\alpha\rVert$.  Since
$N^1(X)_{\mathbf Q}$ is dense in $N^1(X)_{\mathbf R}$, there is a
$\mathbf Q$-line bundle $L$ with $c_1(L)\in U$; the first two
assertions hold by construction.

For bigness, write $\beta:=c_1(L)-\theta$, an ample real class.
Since $\operatorname{Psef}(X)$ is by definition the closure of the convex
cone generated by classes of effective Cartier divisors, there are effective
$\mathbf Q$-Cartier divisors $E_i$ with $c_1(E_i)\to\theta$.  The set
$\{\gamma:\beta+\theta-\gamma\in\Amp(X)\}$ is an open neighbourhood of
$\theta$, so for some $i$ the class
$\beta':=\beta+\theta-c_1(E_i)$ is ample.  This class is rational,
since $\beta'=c_1(L)-c_1(E_i)$ and both $L$ and $E_i$ are rational.  Then
\[
 c_1(L)=c_1(E_i)+\beta' .
\]
Choose $N$ so that $NE_i$ is an integral effective divisor and
$N\beta'$ is the class of an ample line bundle $A'$, and let
$s\in H^0(X,\cO_X(NE_i))$ be the canonical section.  Multiplication by
$s^{\,m}$ is an injection
$H^0(X,mA')\hookrightarrow H^0\bigl(X,m(A'+NE_i)\bigr)$ for every
$m\ge1$, whence
\[
 h^0\bigl(X,m(A'+NE_i)\bigr)\ \ge\ h^0(X,mA')
 \ \sim\ \frac{(A'^{\,n})}{n!}\,m^n,
 \qquad n=\dim X,
\]
by asymptotic Riemann--Roch for the ample bundle $A'$.  Since
$(A'^{\,n})>0$, the line bundle $A'+NE_i$ is big.  Its class is
$N\beta'+Nc_1(E_i)=N\,c_1(L)$, and bigness depends only on the numerical
class \cite[Corollary~2.2.8]{PAG1}, so $c_1(L)$ is big.  (One cannot
shorten this by writing $\cO_X(mNc_1(L))$: $c_1(L)$ is a numerical class,
$A'+NE_i$ is merely numerically equivalent to $NL$, and $h^0$ is not a
numerical invariant.)
\end{proof}

\begin{remark}[Why the drift is oriented and not merely small]
\label{rem:ch40-reduction-orientation}
\Cref{lem:ch40-ep-nef-limit} requires $\theta_j-\theta$ nef, not merely
$\theta_j\to\theta$; \Cref{lem:ch40-rational-approximation} supplies the
stronger property that $\theta_j-\theta$ is \emph{ample}.  This is the
discipline of \Cref{prop:ch38-drift-ledger}\,\cref{itm:ch38-drift:orient} and
\Cref{warn:ch38-drift-not-absorbed}, applied here in the direction that makes
$\PSH_X(\theta)\subseteq\PSH_X(\theta_j)$.  The stage classes in
\Cref{sec:ch40-stages} have the same nef-cone orientation: they too dominate
the target.  What differs is the limiting argument.  Here a family is lifted
to nearby big rational classes, with no monotonicity required of those
classes; there the stage classes form a nef-decreasing sequence while the
stage functions decrease to the target.  The two situations are handled by
different clauses of
\Cref{thm:bj-basic-closure}: clause~(3) here, clauses~(4) and~(5) there.
Keeping them apart is the reason both are displayed.
\end{remark}

\subsection{From families to rational PL obstacles}
\label{sec:ch40-families-to-obstacles}

The passage from families to obstacles is \Cref{thm:general-obstacle-criterion},
proved in \Cref{chap:bj-envelope-history} without any envelope input.  We
recall precisely which implication is used, since the theorem is an
equivalence and one of its directions is forbidden by
\Cref{warn:ch40-firewall}.

\begin{lemma}
\label{lem:ch40-obstacle-interface}
Let $X$ be integral projective and let $\theta\in N^1(X)_{\mathbf R}$
satisfy $\PSH_X(\theta)\ne\varnothing$.  If
$P_\theta(h)\in\PSH_X(\theta)$ for every $h\in\PL(X)$, then
$\operatorname{EP}_X(\theta)$ holds.  Moreover, if
$\theta=c_1(L)$ with $L$ a big $\mathbf Q$-line bundle, then the
hypothesis $\PSH_X(\theta)\ne\varnothing$ is automatic.
\end{lemma}

\begin{proof}
The first assertion is the implication (2)$\Rightarrow$(1) of
\Cref{thm:general-obstacle-criterion}.  For the second, bigness of $L$
gives $m\ge1$ with $mL$ an honest line bundle admitting a nonzero
section $s$; then $m^{-1}\log|s|$ is the Fubini--Study function of the
one-element weighted section system $(m,(s),(0))$, and it is generically
finite because $X$ is integral and $s\ne0$.  Hence
$m^{-1}\log|s|\in\cH^{\mathrm{gf}}(L)\subseteq\PSH_X(\theta)$ by
\Cref{def:bj-psh}, the constant net being decreasing.
\end{proof}

\begin{warning}{Two obstacle envelopes, one symbol}
\label{warn:ch40-two-envelopes}
\Cref{def:bj-obstacle-envelope} defines $P_\theta(h)$ as the \emph{raw}
supremum of the competitors, whereas
\eqref{eq:bj-smooth-obstacle-envelope-definition} writes $P_L(\phi_D)$
for its upper-semicontinuous regularization.  For a \emph{continuous}
obstacle the two agree as soon as either is psh, because a psh function is
usc and lies below the continuous majorant $h$, so the regularization of
the raw supremum is again a competitor; this is the argument in the proof of
\Cref{thm:general-obstacle-criterion} for PL obstacles.  Before psh-ness is
known the two need not agree, and the reader should keep in mind that all
inequalities proved below are proved for the raw supremum and survive
regularization only in the direction in which they are used.  In particular
the inequality $\tau\le u$ of \eqref{eq:bj-smooth-limit-below-obstacle} is
proved for each competitor $\tau$ separately and only then passed to the
supremum, never the other way round.
\end{warning}

\subsection{The interface proposition}
\label{sec:ch40-interface}

\begin{proposition}\label{prop:bj-smooth-obstacle-interface}
To prove $\operatorname{EP}_X(\theta)$ for every real numerical class, it
is enough to prove the following assertion.  If $L$ is a big rational
line bundle and $\phi_D$ is a rational PL obstacle, then the psh obstacle
envelope
\[
       P_L(\phi_D):=
       \left(\sup\{\tau\in\PSH_X(L):\tau\le\phi_D\}\right)^*
\tag{40.1}\label{eq:bj-smooth-obstacle-envelope-definition}
\]
belongs to $\PSH_X(L)$.
\end{proposition}

\begin{proof}
Assume the displayed assertion, and let $\theta\in N^1(X)_{\mathbf R}$ be
arbitrary.

\emph{Step 1: non-pseudoeffective classes.}  If
$\theta\notin\operatorname{Psef}(X)$, then
$\PSH_X(\theta)=\varnothing$ by \Cref{lem:ch34-psh-nonempty}.  There is
then no nonempty family to which
\Cref{def:bj-envelope-property} could apply, so
$\operatorname{EP}_X(\theta)$ holds vacuously.  This is the case recorded
in \Cref{ex:ch37-empty-class}.

\emph{Step 2: reduction to big rational classes.}  Suppose
$\theta\in\operatorname{Psef}(X)$.  By
\Cref{lem:ch40-rational-approximation} applied with $\delta=1/j$ there are
$\mathbf Q$-line bundles $L_j$, $j\ge1$, with $c_1(L_j)$ big,
$c_1(L_j)-\theta$ ample, and $c_1(L_j)\to\theta$.  By
\Cref{lem:ch40-ep-nef-limit} it suffices to prove
$\operatorname{EP}_X(c_1(L_j))$ for every $j$.  We may therefore assume
from now on that $\theta=c_1(L)$ with $L$ a big $\mathbf Q$-line
bundle.

\emph{Step 3: reduction to PL obstacles.}  By
\Cref{lem:ch40-obstacle-interface} it suffices to prove that
$P_L(h)\in\PSH_X(c_1(L))$ for every $h\in\PL(X)$.

\emph{Step 4: normalizing the obstacle.}  Fix $h\in\PL(X)$.  By
\Cref{thm:tc-pl-correspondence} there are a test configuration
$\cY$ for $X$ dominating $X\times\mathbf A^1$ through a morphism
$\mu$ and a vertical $\mathbf Q$-Cartier divisor $D_0$ on $\cY$ with
$h=\phi_{D_0}$; the divisor is vertical because $h$ is the difference of
two flag-ideal functions, whose associated divisors are supported over
$0\in\mathbf A^1$.  Replacing $\cY$ by the normalization of the closure of
$X\times\mathbf G_m$ inside it, and $D_0$ by its pullback, we may and do
assume that $\cY$ is integral and normal; this changes neither $h$ nor
the class of the test configuration, by the pullback invariance
$\phi_{\mu^*D}=\phi_D$ established in Step~1 of the proof of
\Cref{thm:tc-pl-correspondence}.  Being vertical, $D_0$ is bounded above by
$N\cY_0$ for some $N\in\mathbf Q_{>0}$, where $\cY_0$ is the central
fibre; replacing $D_0$ by $D_1:=D_0-N\cY_0$ replaces $h$ by $h-N$,
by \eqref{eq:ch34-flag-weight-form} and the normalization
$\phi_{\cY_0}=1$ of the convention sheet, and
$P_L(h-N)=P_L(h)-N$ by \Cref{lem:obstacle-envelope-formal}\,(2).  Since
$\PSH_X(c_1(L))$ is stable under adding real constants,
\Cref{thm:bj-basic-closure}\,(1), we may assume $D_0\le0$.

Finally, choose $r\in\mathbf Z_{>0}$ so that $rL$ is an honest line
bundle and $rD_0$ is an integral Cartier divisor, and put
$L':=rL$, $D:=rD_0$, $h':=rh=\phi_D$.  Scaling by $r$ is a bijection
$\PSH_X(c_1(L))\to\PSH_X(c_1(L'))$, $\psi\mapsto r\psi$, by
\Cref{thm:bj-basic-closure}\,(3) applied with $s=r,t=0$ and its inverse
with $s=r^{-1}$; it carries competitors for $h$ to competitors for
$h'$ and therefore satisfies $P_{L'}(h')=rP_L(h)$, regularization
commuting with multiplication by the positive constant $r$.  Hence
$P_L(h)$ is $c_1(L)$-psh if and only if $P_{L'}(\phi_D)$ is
$c_1(L')$-psh.  As $L'$ is again big, this is the displayed assertion
for the rational PL obstacle $\phi_D$ with $D\le0$, and the proof is
complete.
\end{proof}

\begin{remark}[What the reduction costs, and what it does not]
\label{rem:ch40-reduction-cost}
Three features of the reduction should be recorded, because each is a place
where a plausible variant fails.

\emph{The four steps do not commute.}  Step~2 must precede Step~3: the
obstacle criterion \Cref{thm:general-obstacle-criterion} needs
$\PSH_X(\theta)\ne\varnothing$, and for a merely pseudoeffective real
$\theta$ that is not automatic in a usable form, whereas for a big
rational class it is \Cref{lem:ch40-obstacle-interface}.  Conversely Step~4
must follow Step~3, because the normalization $D\le0$ refers to a
particular obstacle and there is no ``uniformly normalized'' family of
obstacles to work with.

\emph{Bigness, not ampleness.}  The reduction produces a \emph{big} class
and no more.  This matters: $\PSH_X(c_1(L))$ then contains functions
identically $-\infty$ on the augmented base locus, the base ideals
$\mathfrak a_m$ of \Cref{sec:ch40-base-ideals} have nonempty zero locus,
and the final comparison in
\Cref{thm:bj-smooth-envelope} has to be made at divisorial points precisely
in order to avoid that locus.  A reader who imagines $L$ ample will find
several of the extended-real precautions below unmotivated;
\Cref{ex:bj-common-base-locus} is the example to keep in view.

\emph{Nothing here is characteristic-dependent.}  Steps~1--4 use only
\Cref{thm:bj-basic-closure}, the cone geometry of $N^1(X)_{\mathbf R}$,
and \Cref{thm:tc-pl-correspondence}.  The characteristic enters for the
first time in \Cref{sec:ch40-engines}.  The reduction is also valid for
singular $X$; it is \Cref{sec:ch40-engines} onwards that needs $X$
smooth, and \Cref{sec:ch40-audit} records exactly where.
\end{remark}

\section{Graded base ideals of an obstacle}
\label{sec:ch40-base-ideals}

By \Cref{prop:bj-smooth-obstacle-interface} we are reduced to the following
situation, which is fixed for the rest of the chapter.  After clearing
denominators, taking a Veronese, and adding a constant to the obstacle, its
normalized graph over the trivial test configuration has
the form
\[
 \mu:\mathcal Y\longrightarrow
 S:=X\times\mathbf A^1,\qquad
 \mathcal L=\mu^*L+D,\qquad D\le0.
\tag{40.2}\label{eq:bj-smooth-graph}
\]
Here $L$ is a big line bundle on $X$, $\mu$ is an equivariant
projective birational morphism from an integral normal scheme $\cY$
carrying a $\mathbf G_m$-action lifting the scaling action on $S$, and
$D$ is a $\mathbf G_m$-invariant integral Cartier divisor supported over
$0\in\mathbf A^1$.  The obstacle is $\phi_D$, and the objective is
\[
 P_L(\phi_D)\in\PSH_X(c_1(L)).
\]
Write
\[
 \rho=\operatorname{pr}_{\mathbf A^1}:S\to\mathbf A^1,\qquad
 \pi=\rho\circ\mu,\qquad
 M_m=\pi_*\cO_{\mathcal Y}(m\mathcal L).
\]
The $k[\varpi]$-module $M_m$ is coherent and torsion free: a section killed
by a nonzero polynomial vanishes on the generic fibre and hence on the
integral scheme $\mathcal Y$.  Since $k[\varpi]$ is a PID, $M_m$ is finite
free.

Two comments on this normalization.  First, $\cY$ may and will be taken
\emph{normal} but is not assumed smooth; no resolution of $\cY$ is
performed anywhere in this chapter, and this is one of the two structural
differences from the argument of \Cref{thm:bj22-nef-envelope}, which resolves
$\cX$ and consequently needs $\dim\cX\le3$ in positive characteristic.
Second, all the ideal theory below takes place on $S=X\times\mathbf A^1$
and not on $\cY$: $S$ is smooth over $k$, hence regular and (in
characteristic $p$) $F$-finite, whereas $\cY$ is merely normal.  The
morphism $\mu$ is used only to \emph{define} the ideals and to prove the
identity \eqref{eq:bj-smooth-obstacle-identity}; it never carries an ideal
sheaf whose multiplier or test ideal is taken.

\subsection{The ambient graded sequence}
\label{sec:ch40-ambient-ideal}

Let $\mathfrak a_m$ be the relative base ideal of $m\mathcal L$ on
$\mathcal Y$.  Since $S$ is normal and $\mu$ is proper birational,
$\mu_*\cO_{\mathcal Y}=\cO_S$.  Define an equivariant coherent ideal
$J_m\subset\cO_S$ by
\[
 M_m\otimes\cO_S(-mL)\longrightarrow
 \mu_*\cO_{\mathcal Y}(mD)\subset\cO_S,\qquad
 J_m:=\operatorname{im}.
\tag{40.3}\label{eq:bj-smooth-ambient-ideal}
\]
In \eqref{eq:bj-smooth-ambient-ideal} the module $M_m$ is pulled back to
$S$ along $\rho$, and the arrow is the evaluation map of the relative
linear system, twisted down by $mL$; its target is contained in
$\cO_S$ because $D\le0$.
After the chosen Veronese, $J_m\ne0$ in every degree, and multiplication
of sections gives
\[
                            J_mJ_r\subset J_{m+r}.
\tag{40.4}\label{eq:bj-smooth-graded-ideals}
\]
Moreover, $J_q(qL)$, and hence $J_q^r(qrL)$, is generated relative to
$\mathbf A^1$.

We record these assertions as a lemma and prove them, since every later
construction rests on all four.

\begin{lemma}
\label{lem:ch40-J-basic}
In the situation \eqref{eq:bj-smooth-graph}:
\begin{enumerate}
\item\label{itm:ch40-J:coherent}
 $J_m\subseteq\cO_S$ is a coherent $\mathbf G_m$-equivariant ideal, and
 $J_m(mL)$ is generated relative to $\mathbf A^1$; consequently so is
 $J_m^{\,r}(mrL)$ for every $r\ge1$;
\item\label{itm:ch40-J:nonzero}
 there is $m_1\ge1$ such that $J_m\ne0$ for every $m\in m_1\mathbf Z_{>0}$;
\item\label{itm:ch40-J:graded}
 $J_mJ_r\subseteq J_{m+r}$ for all $m,r\ge1$;
\item\label{itm:ch40-J:pullback}
 $J_m\cO_{\cY}=\mathfrak a_m\cO_{\cY}(mD)$.
\end{enumerate}
\end{lemma}

\begin{proof}
\cref{itm:ch40-J:coherent}  The map in \eqref{eq:bj-smooth-ambient-ideal} is
a map of coherent sheaves, so its image is coherent; it is equivariant
because $\cL$, $D$ and $\mu$ are.  Its target is
$\mu_*\cO_{\cY}(mD)$, which is a subsheaf of
$\mu_*\cO_{\cY}=\cO_S$ because $mD\le0$; hence $J_m\subseteq\cO_S$.
Twisting \eqref{eq:bj-smooth-ambient-ideal} by $mL$ exhibits
$J_m(mL)$ as the image of $\rho^*M_m=\rho^*\rho_*\bigl(\mu_*\cO_{\cY}(m\cL)\bigr)$,
which is precisely the assertion that $J_m(mL)$ is generated relative to
$\mathbf A^1$.  If $\mathcal F$ and $\mathcal G$ are relatively
generated then so is $\mathcal F\otimes\mathcal G$, and the image of
$\mathcal F\otimes\mathcal G$ under a surjection is relatively generated;
applying this $r$ times to $J_m(mL)$ gives the last claim, since the
multiplication map $J_m(mL)^{\otimes r}\to J_m^{\,r}(mrL)$ is surjective.

\cref{itm:ch40-J:nonzero}  Restricting \eqref{eq:bj-smooth-ambient-ideal} to
the open set $\{\varpi\ne0\}\cong X\times\mathbf G_m$, where $\mu$ is an
isomorphism and $\cL$ is the pullback of $L$, identifies $J_m$ with the
base ideal of the linear system $|mL|$ on $X$, pulled back.  Since $L$
is big there is $m_1$ with $H^0(X,m_1L)\ne0$, and then
$H^0(X,mL)\ne0$, hence $J_m\ne0$, for every positive multiple $m$ of
$m_1$.  We replace $L$ and $D$ by $m_1L$ and $m_1D$ once and for
all, so that $J_m\ne0$ for every $m\ge1$; this is the Veronese referred
to in \eqref{eq:bj-smooth-ambient-ideal}, and it changes neither the
obstacle problem, by the scaling bijection used in Step~4 of the proof of
\Cref{prop:bj-smooth-obstacle-interface}, nor any of the functions below.

\cref{itm:ch40-J:graded}  If $s\in M_m$ and $s'\in M_r$, then
$ss'\in M_{m+r}$, and under \eqref{eq:bj-smooth-ambient-ideal} the germ of
$ss'$ is the product of the germs of $s$ and $s'$.  Since $J_m$ is
generated by the germs of sections in $M_m$, the product ideal
$J_mJ_r$ is generated by such products and is therefore contained in
$J_{m+r}$.

\cref{itm:ch40-J:pullback}  Pull \eqref{eq:bj-smooth-ambient-ideal} back
along $\mu$.  Since $\mu^*$ is right exact, the image of
$\mu^*\bigl(M_m\otimes\cO_S(-mL)\bigr)$ in $\cO_{\cY}$ is
$J_m\cO_{\cY}$.  On the other hand the same map factors as
\[
 \mu^*\rho^*M_m\otimes\cO_{\cY}(-m\mu^*L)
 \longrightarrow\cO_{\cY}(mD)
 \hookrightarrow\cO_{\cY},
\]
whose image is by definition $\mathfrak a_m\cO_{\cY}(mD)$, the relative
base ideal of $m\cL=m\mu^*L+mD$ twisted back down.  Comparing the two
descriptions gives \cref{itm:ch40-J:pullback}.
\end{proof}

Pulling \eqref{eq:bj-smooth-ambient-ideal} to $\mathcal Y$ and using
right exactness gives
\[
 J_m\cO_{\mathcal Y}
   =\mathfrak a_m\cO_{\mathcal Y}(mD).
\tag{40.5}\label{eq:bj-smooth-base-ideal-equality}
\]
Read this as an equality of ideal values before changing signs.  If
$\phi_I(v)=-\sigma(v)(I)$, where the Gauss extension is normalized by
$\sigma(v)(\varpi)=1$, then
\[
 \frac1m\phi_{J_m}
   =\phi_D+\frac1m\phi_{\mathfrak a_m}.
\tag{40.6}\label{eq:bj-smooth-obstacle-identity}
\]
This identity remains literal when one ideal value is $+\infty$.

\begin{proof}[Proof of the obstacle identity]
Fix $v\in X^{\An}$ and let $w:=\sigma(v)$ be its Gauss extension, a
semivaluation on $\cO_S$ with $w(\varpi)=1$.  Since $\mu$ is proper
and surjective, $w$ admits a lift $\tilde w$ to $\cO_{\cY}$,
that is, a semivaluation with $\tilde w\circ\mu^\sharp=w$.  Indeed, write
$w=(\xi,|\cdot|)$ as in \Cref{lem:ch41-point-pair}, choose a point
$\eta\in\cY$ with $\mu(\eta)=\xi$, which exists because $\mu$ is
surjective, and extend $|\cdot|$ from $\kappa(\xi)$ to the field
extension $\kappa(\eta)$; an absolute value extends to an arbitrary field
extension by \cite[Chapter~VI, \S4]{ZariskiSamuelII}.  For an ideal
$I\subseteq\cO_S$ one then has
$\tilde w(I\cO_{\cY})=w(I)$, because a finite set of local generators of
$I$ generates $I\cO_{\cY}$ and the value of an ideal is the minimum of
the values of a generating set, \eqref{eq:ideal-evaluation-triv}.  Applying
this to \eqref{eq:bj-smooth-base-ideal-equality} and using
$\tilde w(\mathcal I\mathcal J)=\tilde w(\mathcal I)+\tilde w(\mathcal J)$,
\Cref{prop:ch32-tropical-identities}\,\cref{itm:ch32-ti:identities}, gives
\[
 w(J_m)=\tilde w(\mathfrak a_m)+\tilde w\bigl(\cO_{\cY}(mD)\bigr).
\]
Now $\cO_{\cY}(mD)$ is the ideal sheaf of the effective divisor
$-mD$, so $\tilde w(\cO_{\cY}(mD))=-m\phi_D(v)$ by the sign convention of
the front matter, recalled in \Cref{warn:ch40-signs}, and
$\tilde w(\mathfrak a_m)=-\phi_{\mathfrak a_m}(v)$.  Negating and dividing
by $m$ gives \eqref{eq:bj-smooth-obstacle-identity}.  All three terms lie
in $[0,+\infty]$ before negation, and the identity is an identity in
$[0,+\infty]$; no difference of two infinite values is formed, in
accordance with the convention sheet.
\end{proof}

\begin{warning}{The two sign tests}
\label{warn:ch40-signs}
Every sign in this chapter is fixed by two unit tests, and a reader who
checks them will not need to re-derive any other.

\emph{Test one: an effective vertical divisor gives a nonnegative function.}
If $E$ is an effective vertical Cartier divisor on $\cY$ with local
equation $f$, then $\cO_{\cY}(E)$ is generated by $f^{-1}$ and
$\phi_E(v)=\tilde w(f)\ge0$.  Hence $D\le0$ gives $\phi_D\le0$, which
is why the normalization of Step~4 in
\Cref{prop:bj-smooth-obstacle-interface} was arranged.  In particular
$\phi_{\cY_0}=1$ for the central fibre.

\emph{Test two: an ideal gives a nonpositive function.}  For
$I\subseteq\cO_S$, $\phi_I=-\sigma(\cdot)(I)\le0$, with equality exactly
where $I$ is locally the unit ideal at the centre.  Thus larger ideals give
larger functions:
$I\subseteq I'\Rightarrow\phi_I\le\phi_{I'}$.  This order reversal is the
mechanism that turns the subadditivity $B_{m\ell}\subseteq B_m^{\,\ell}$
into the monotonicity $u_{m\ell}\le u_m$, and it is the single most
important sign in the chapter.

Combining the two tests: in \eqref{eq:bj-smooth-obstacle-identity} the
left-hand side is $\le\phi_D$ because
$\mathfrak a_m\subseteq\cO_{\cY}$, and it increases with $m$ in the
divisibility order because $\mathfrak a_m\mathfrak a_r
\subseteq\mathfrak a_{m+r}$.
\end{warning}

\subsection{Equivariant ideals as Fubini--Study functions}
\label{sec:ch40-equivariant-fs}

The bridge between the ideal-theoretic constructions and the psh formalism
is the following lemma.  It is the exact analogue of
\Cref{prop:ch34-fs-flag-ideal} for equivariant ideals that are \emph{not}
trivial over $\mathbf G_m$; the difference matters, because $L$ is only
big, so the base ideal of $|mL|$ has nonempty zero locus and the ideals
$J_m$ and $B_m$ below are not flag ideals in the sense of
\Cref{def:bj-flag-ideal}.

\begin{lemma}
\label{lem:ch40-equivariant-ideal-fs}
Let $N$ be a $\mathbf Q$-line bundle on $X$, let $m\ge1$ be such
that $mN$ is an honest line bundle, and let $0\ne I\subseteq\cO_S$ be a
coherent $\mathbf G_m$-equivariant ideal.
\begin{enumerate}
\item\label{itm:ch40-eqfs:decomp}
 There is a unique increasing sequence of coherent ideals
 $\mathfrak b_0\subseteq\mathfrak b_1\subseteq\cdots\subseteq\cO_X$,
 stationary for $k\gg0$, with
 $I=\bigoplus_{k\ge0}\mathfrak b_k\varpi^k$ as a sheaf of
 $\cO_X[\varpi]$-modules, and for every $v\in X^{\An}$
 \[
   \sigma(v)(I)=\min_{k\ge0}\{v(\mathfrak b_k)+k\},
   \qquad
   \phi_I=\max_{k\ge0}\{\log|\mathfrak b_k|-k\},
 \]
 a finite maximum.
\item\label{itm:ch40-eqfs:gg}
 $I(mN)$ is generated relative to $\mathbf A^1$ if and only if
 $\mathfrak b_k(mN)$ is globally generated on $X$ for every $k$.
\item\label{itm:ch40-eqfs:fs}
 If $I(mN)$ is generated relative to $\mathbf A^1$, then
 \[
   \frac1m\phi_I\in\cH^{\mathrm{gf}}_{m^{-1}\mathbf Z}(N)
   \subseteq\cH^{\mathrm{gf}}_{\mathbf Q}(N),
 \]
 with a presentation of degree $m$ whose weights are the rational
 numbers $-k/m$, for the integers $k$ occurring in
 \cref{itm:ch40-eqfs:decomp}.
\end{enumerate}
\end{lemma}

\begin{proof}
\cref{itm:ch40-eqfs:decomp}  Work on an affine chart
$\Spec A\subseteq X$, so that the chart of $S$ is
$\Spec A[\varpi]$ with the $\mathbf G_m$-action given by the grading in
which $A$ has weight $0$ and $\varpi$ has weight $1$.  By
\Cref{lem:ch33-weight-decomposition} an equivariant submodule is the direct
sum of its weight pieces, so $I=\bigoplus_k\mathfrak b_k\varpi^k$ with
$\mathfrak b_k\subseteq A$ an ideal.  Multiplication by $\varpi$ maps
$I$ into $I$ and raises the weight by one, whence
$\mathfrak b_k\subseteq\mathfrak b_{k+1}$; the chain is stationary because
$A$ is noetherian.  The pieces are independent of the chart, being weight
spaces, so they glue.

For the value of $I$, choose local generators of $\mathfrak b_k$ for the
finitely many $k\le k_0$ up to which the chain is not yet stationary; the
elements $g\varpi^k$, for $g$ such a generator, generate $I$.  By
\eqref{eq:ch40-gauss}, $\sigma(v)(g\varpi^k)=v(g)+k$, and by
\eqref{eq:ideal-evaluation-triv} the value of an ideal is the minimum over a
generating set; hence
$\sigma(v)(I)=\min_{k\le k_0}\{v(\mathfrak b_k)+k\}$.  For $k>k_0$ one
has $\mathfrak b_k=\mathfrak b_{k_0}$ and $k>k_0$, so those terms are
larger and the minimum over all $k\ge0$ is the same.  Negating gives the
formula for $\phi_I$.

\cref{itm:ch40-eqfs:gg}  Both $\rho^*\rho_*(I(mN))$ and $I(mN)$ are
$\mathbf Z$-graded for the $\varpi$-grading, and the evaluation map is
homogeneous of degree $0$; a map of graded sheaves is surjective if and
only if it is surjective in each degree.  In degree $k$ the source is
$H^0(X,\mathfrak b_k(mN))\otimes_k\cO_X$ and the target is
$\mathfrak b_k(mN)$, because $\rho_*(I(mN))
=\bigoplus_kH^0(X,\mathfrak b_k(mN))\varpi^k$ as a $k[\varpi]$-module.
The assertion follows.

\cref{itm:ch40-eqfs:fs}  By \cref{itm:ch40-eqfs:gg} choose, for each
$k\le k_0$, sections $s_{k,1},\ldots,s_{k,r_k}\in H^0(X,mN)$ generating
$\mathfrak b_k(mN)$.  Computing in a frame of $mN$ at the centre of
$v$, these germs generate $(\mathfrak b_k)_{c(v)}$, so
$\log|\mathfrak b_k|=\max_i\log|s_{k,i}|$ as in the proof of
\Cref{prop:ch34-fs-flag-ideal}.  Therefore
\[
 \frac1m\phi_I
 =\frac1m\max_{k\le k_0}\max_{1\le i\le r_k}
  \bigl\{\log|s_{k,i}|+(-k)\bigr\},
\]
a weighted section system of degree $m$ for $N$ with integer weights
$-k$.  It is generically finite: since $I\ne0$ and $X$ is integral,
some $\mathfrak b_k\ne0$, so some $s_{k,i}$ is a nonzero section and the
common zero locus of the whole system is a proper closed subset; now apply
\Cref{prop:ch34-generic-finiteness}\,\cref{itm:ch34-gf:equiv}.
\end{proof}

\begin{remark}[Why relative generation, and not global generation]
\label{rem:ch40-relative-generation}
\Cref{lem:ch40-equivariant-ideal-fs} asks for generation relative to
$\mathbf A^1$, which for the affine base $\mathbf A^1$ is the same as
global generation of the $k[\varpi]$-module $\rho_*(I(mN))$ and, by
part~\cref{itm:ch40-eqfs:gg}, the same as global generation of the finitely
many ideals $\mathfrak b_k(mN)$ on $X$.  The relative formulation is the
one that survives the vanishing arguments of \Cref{sec:ch40-stages}, where
the higher direct images are taken along $\rho$; the fibrewise formulation
is the one that produces the sections.  The equivalence of the two is
exactly part~\cref{itm:ch40-eqfs:gg}, and it is used in both directions.
\end{remark}

\subsection{The two gradings, on the ideals of this chapter}
\label{sec:ch40-two-gradings}

\Cref{prop:ch38-two-gradings} distinguished the \emph{degree} $m$ of a
weighted section system from the \emph{weight} $\lambda$.  In the present
chapter both indices are in play simultaneously and on the same object, so
it is worth saying explicitly what each of them is and how each behaves.

\begin{proposition}
\label{prop:ch40-two-gradings-here}
Let $J_m=\bigoplus_{k\ge0}\mathfrak b_{m,k}\varpi^k$ be the weight
decomposition of \Cref{lem:ch40-equivariant-ideal-fs}\,\cref{itm:ch40-eqfs:decomp}.
\begin{enumerate}
\item\label{itm:ch40-grad:degree}
 The index $m$ is the \emph{degree}: it is the exponent of $\cL$, the
 twist by which $J_m(mL)$ is relatively generated, and the denominator in
 $q_m=m^{-1}\phi_{J_m}$.  Under $m\mapsto rm$ it behaves
 multiplicatively:
 $J_m^{\,r}\subseteq J_{rm}$, and
 $\mathfrak b_{m,k}^{\,r}\subseteq\mathfrak b_{rm,rk}$.
\item\label{itm:ch40-grad:weight}
 The index $k$ is the $\mathbf G_m$-\emph{weight}: it is the exponent of
 $\varpi$, the filtration index of \Cref{sec:ch33-norms}, and it enters
 the Fubini--Study presentation as the additive weight $-k/m$.  Under
 $m\mapsto rm$ the relevant weights are multiplied by $r$ as well, so the
 \emph{ratio} $k/m$ is the invariant quantity.
\item\label{itm:ch40-grad:errors}
 An error of size $O(m^{-1})$ in the first index is a perturbation of the
 numerical class by $O(m^{-1})$; an error of size $O(1)$ in the second
 index is a perturbation of the function by the additive constant
 $O(m^{-1})$.  The two are corrected by different devices and never cancel
 each other.
\item\label{itm:ch40-grad:homogeneous}
 The function $q_m$ is homogeneous, in the sense of
 \Cref{def:homogeneous-function}, if and only if
 $\mathfrak b_{m,k}\subseteq\overline{\mathfrak b_{m,0}}$ for every
 $k\ge0$, the bar denoting integral closure; in that case
 $q_m=m^{-1}\log|\mathfrak b_{m,0}|$.  In particular this holds when
 $J_m=\mathfrak b_{m,0}\cO_S$ is extended from $X$, and it fails
 whenever $\mathfrak b_{m,0}=0$.
\end{enumerate}
\end{proposition}

\begin{proof}
\cref{itm:ch40-grad:degree}  The inclusion $J_m^{\,r}\subseteq J_{rm}$ is
\Cref{lem:ch40-J-basic}\,\cref{itm:ch40-J:graded} applied $r-1$ times.
Taking weight-$rk$ pieces of both sides gives the second inclusion,
because the weight-$rk$ piece of $J_m^{\,r}$ contains
$(\mathfrak b_{m,k}\varpi^k)^r=\mathfrak b_{m,k}^{\,r}\varpi^{rk}$.  The
remaining assertions are the definitions.

\cref{itm:ch40-grad:weight}  The identification of $k$ with the filtration
index is \Cref{lem:ch33-flag-weight-decomposition} and
\Cref{rem:ch33-two-indexings}; the appearance of $-k/m$ as a weight is
\Cref{lem:ch40-equivariant-ideal-fs}\,\cref{itm:ch40-eqfs:fs}.  If
$m$ is replaced by $rm$ then, by \cref{itm:ch40-grad:degree}, the weight
that carries the same information is $rk$, and
$-rk/(rm)=-k/m$.

\cref{itm:ch40-grad:errors}  This is
\Cref{warn:ch38-two-gradings} transcribed to the present notation, together
with \Cref{rem:ch38-two-errors}.  \Cref{sec:ch40-stages} exhibits the first
error, of exact size $m^{-1}c_1(A)$; the second does not occur below,
because every weight used is an honest integer divided by the honest degree.

\cref{itm:ch40-grad:homogeneous}  By
\Cref{lem:ch40-equivariant-ideal-fs}\,\cref{itm:ch40-eqfs:decomp},
$q_m=m^{-1}\max_k\{\log|\mathfrak b_{m,k}|-k\}$.

Suppose first that $\mathfrak b_{m,k}\subseteq\overline{\mathfrak b_{m,0}}$
for every $k$.  Ideal values do not distinguish an ideal from its integral
closure, by \Cref{cor:appA-integral-closure-values}, so
$\log|\mathfrak b_{m,k}|\le\log|\overline{\mathfrak b_{m,0}}|
=\log|\mathfrak b_{m,0}|$; since moreover $-k\le0$, the maximum is
attained at $k=0$ and $q_m=m^{-1}\log|\mathfrak b_{m,0}|$, which is
homogeneous by \Cref{prop:hom-fs-ideal}.

Conversely, suppose $q_m$ is homogeneous.  If $\mathfrak b_{m,0}=0$ then
the presentation of
\Cref{lem:ch40-equivariant-ideal-fs}\,\cref{itm:ch40-eqfs:fs} has all its
weights $-k/m$ strictly negative, so its largest weight is nonzero,
contradicting
\Cref{prop:ch38-two-gradings}\,\cref{itm:ch38-grad:weight}; hence
$\mathfrak b_{m,0}\ne0$.  Let $k_0$ be the index at which the increasing
chain $(\mathfrak b_{m,k})_k$ becomes stationary, so that
$\log|\mathfrak b_{m,k}|\le\log|\mathfrak b_{m,k_0}|$ for every $k$.
The same item asserts that the weights may be deleted from any presentation
without changing the function, which for the presentation of
\Cref{lem:ch40-equivariant-ideal-fs}\,\cref{itm:ch40-eqfs:fs} reads
\[
 \max_k\bigl\{\log|\mathfrak b_{m,k}|-k\bigr\}
 =\max_k\log|\mathfrak b_{m,k}|
 =\log|\mathfrak b_{m,k_0}| .
\]
Fix $v\in X^{\val}$, where all these values are finite.  The left-hand
maximum is attained at some $k$, and
$\log|\mathfrak b_{m,k}|(v)-k\le\log|\mathfrak b_{m,k_0}|(v)-k$; equality
with $\log|\mathfrak b_{m,k_0}|(v)$ therefore forces $k=0$ and
$\log|\mathfrak b_{m,0}|(v)=\log|\mathfrak b_{m,k_0}|(v)$.  As $v$ was
arbitrary and the chain is increasing, $\log|\mathfrak b_{m,k}|
=\log|\mathfrak b_{m,0}|$ on $X^{\val}$ for every $k$.  Evaluating at
divisorial valuations gives $v(\mathfrak b_{m,k})=v(\mathfrak b_{m,0})$,
which by the valuative criterion of integral dependence
\Cref{thm:appA-valuative-integral-closure} is exactly
$\overline{\mathfrak b_{m,k}}=\overline{\mathfrak b_{m,0}}$, and in
particular $\mathfrak b_{m,k}\subseteq\overline{\mathfrak b_{m,0}}$.
\end{proof}

\begin{warning}{The obstacle is not homogeneous}
\label{warn:ch40-two-gradings}
\Cref{chap:bj-reflection} could afford to work with homogeneous functions
throughout, and in that category the second grading is identically zero, by
\Cref{prop:ch38-two-gradings}\,\cref{itm:ch38-grad:weight}.  Here it cannot
be arranged.  The obstacle $\phi_D$ is a PL function attached to a
vertical divisor; it is homogeneous only if $D$ is supported on the
central fibre with the trivial multiplicity structure, which excludes every
interesting case, and \Cref{ex:ch40-p2-obstacle} exhibits a $J_m$ whose
weight decomposition has $m+1$ distinct terms.  A reader who imports the
homogeneous simplifications of \Cref{chap:bj-reflection} into this chapter
will find that the ideals $J_m$ collapse to their weight-zero pieces and
that the resulting ``proof'' computes the envelope of the wrong obstacle,
namely of $\max\{\phi_D,\ \text{a constant}\}$.
\end{warning}

\subsection{The complete base-ideal limit}
\label{sec:ch40-Q}

The left side of \eqref{eq:bj-smooth-obstacle-identity} increases in the
divisibility order to the complete base-ideal limit
\[
                            Q:=Q_L(\phi_D).
\tag{40.7}\label{eq:bj-smooth-Q}
\]
Here $Q_L$ is the Fubini--Study envelope \eqref{eq:fs-envelope}, that is,
the pointwise supremum of all generically finite Fubini--Study functions for
$L$ lying below the obstacle.  We prove the assertion.

\begin{proposition}
\label{prop:ch40-Q-limit}
Set $q_m:=m^{-1}\phi_{J_m}$.  Then:
\begin{enumerate}
\item\label{itm:ch40-Q:gf}
 $q_m\in\cH^{\mathrm{gf}}_{\mathbf Q}(L)$ and $q_m\le\phi_D$ for every
 $m\ge1$;
\item\label{itm:ch40-Q:increasing}
 $q_m\le q_{rm}$ for every $r\ge1$; thus $(q_m)$ is an increasing net
 for the divisibility order;
\item\label{itm:ch40-Q:sup}
 $\sup_mq_m=\lim_mq_m=Q_L(\phi_D)$ pointwise on $X^{\An}$, the limit
 being taken along divisibility.
\end{enumerate}
\end{proposition}

\begin{proof}
\cref{itm:ch40-Q:gf}  $J_m$ is a nonzero coherent equivariant ideal and
$J_m(mL)$ is generated relative to $\mathbf A^1$, by
\Cref{lem:ch40-J-basic}\,\cref{itm:ch40-J:coherent} and
\cref{itm:ch40-J:nonzero}; now apply
\Cref{lem:ch40-equivariant-ideal-fs}\,\cref{itm:ch40-eqfs:fs} with $N=L$.
The inequality $q_m\le\phi_D$ is
\eqref{eq:bj-smooth-obstacle-identity} together with
$\mathfrak a_m\subseteq\cO_{\cY}$, which gives
$\phi_{\mathfrak a_m}\le0$ by the second sign test of
\Cref{warn:ch40-signs}.

\cref{itm:ch40-Q:increasing}  By
\Cref{lem:ch40-J-basic}\,\cref{itm:ch40-J:graded},
$J_m^{\,r}\subseteq J_{rm}$; by the order reversal of
\Cref{warn:ch40-signs} and the identity
$\phi_{I^r}=r\phi_I$ of \eqref{eq:ch34-log-ideal-semiring} this gives
$r\phi_{J_m}\le\phi_{J_{rm}}$, that is, $q_m\le q_{rm}$.

\cref{itm:ch40-Q:sup}  Since each $q_m$ is a generically finite
Fubini--Study function for $L$ lying below $\phi_D$, we have
$q_m\le Q_L(\phi_D)$ and hence $\sup_mq_m\le Q_L(\phi_D)$.

For the reverse inequality, let $\psi\in\cH^{\mathrm{gf}}_{\mathbf R}(L)$
with $\psi\le\phi_D$; we must find $m$ with $\psi\le q_m$, at least
after an arbitrarily small decrease of $\psi$.

\emph{Reduction to integer weights.}  Write
$\psi=m_0^{-1}\max_j\{\log|s_j|+\lambda_j\}$ with $\lambda_j\in\mathbf R$.
For $\varepsilon>0$ choose $\lambda_j'\in\mathbf Q$ with
$\lambda_j-\varepsilon\le\lambda_j'\le\lambda_j$ and put
$\psi_\varepsilon:=m_0^{-1}\max_j\{\log|s_j|+\lambda_j'\}$.  Then
$\psi-\varepsilon/m_0\le\psi_\varepsilon\le\psi\le\phi_D$, and
$\psi_\varepsilon$ has rational weights.  By \Cref{lem:ch34-veronese} we
may pass to a Veronese presentation of degree $m:=rm_0$ with integer
weights $r\lambda_j'$ and sections $s_j^{\,r}\in H^0(X,mL)$.  It
therefore suffices to prove the claim for $\psi$ with integer weights.

\emph{The membership test.}  So let
$\psi=m^{-1}\max_j\{\log|s_j|+\lambda_j\}\le\phi_D$ with
$s_j\in H^0(X,mL)\setminus\{0\}$ and $\lambda_j\in\mathbf Z$.  Fix
$j$ and consider the rational section
$t_j:=s_j\varpi^{-\lambda_j}$ of $m\cL$ on $\cY$; it is regular over
$\mathbf G_m$, where $\varpi$ is invertible and $\cL$ is the pullback
of $L$.  We claim $t_j\in H^0(\cY,m\cL)=M_m$.  Since $\cY$ is normal,
it suffices to show that the divisor of $t_j$ as a section of $m\cL$ is
effective, and for that it suffices to check nonnegativity of
$\ord_P(t_j)$ for every prime divisor $P$ of $\cY$.  If $P$ is
horizontal, that is, dominates $\mathbf A^1$, then $\ord_P(t_j)$ is the
order of vanishing of $s_j$ along the corresponding divisor of $X$,
which is $\ge0$.  If $P=E$ is a component of the central fibre, note first that $\cY$
is normal, hence integrally closed in the sense of
\Cref{def:integrally-closed-tc} by
\Cref{thm:ch33-integrally-closed-criterion}, so that
\Cref{def:ch33-component-valuation} applies; then by
\Cref{lem:ch33-gauss-component} one has $\ord_E=b_E\,\sigma(v_E)$ on
$\cO_S$, with $b_E=\ord_E(\varpi)$ as in \eqref{eq:ch33-bE} and
$v_E\in X^{\An}$; computing in a frame of $m\cL$ obtained as the product
of a frame of $m\mu^*L$ and the local equation of the effective divisor
$-mD$, and using the two sign tests of \Cref{warn:ch40-signs}, we get
\[
 b_E^{-1}\ord_E(t_j)
 =v_E(s_j)-\lambda_j+m\,\phi_D(v_E)
 =m\phi_D(v_E)-\bigl(\log|s_j|(v_E)+\lambda_j\bigr)\ \ge\ 0,
\]
the last inequality being the hypothesis $\psi\le\phi_D$ evaluated at
$v_E$.  Hence $t_j\in M_m$.

\emph{Conclusion.}  The image of $t_j$ under
\eqref{eq:bj-smooth-ambient-ideal} is a local section $g_j$ of $J_m$
with $\sigma(v)(g_j)=v(s_j)-\lambda_j$ for every $v$, computed in the
same frames.  Therefore
$\phi_{J_m}\ge\log|s_j|+\lambda_j$ for every $j$, whence
$q_m\ge\psi$.  Combining with the reduction step,
$\sup_mq_m\ge\psi-\varepsilon/m_0$ for every $\varepsilon>0$, and then
$\sup_mq_m\ge\psi$.  Taking the supremum over $\psi$ gives
$\sup_mq_m\ge Q_L(\phi_D)$.

Finally the supremum is a limit because the net is increasing, by
\cref{itm:ch40-Q:increasing}.
\end{proof}

The relation between $Q$ and the psh obstacle envelope is the following.
Only part~\cref{itm:ch40-cmp:divisorial} is used later; it is proved here in
full, and the reader should note that it is proved for each competitor
separately, in accordance with \Cref{warn:ch40-two-envelopes}.

\begin{proposition}
\label{prop:ch40-envelope-comparison}
Let $L$ be a big line bundle on the integral projective variety $X$, and
let $h:X^{\An}\to\mathbf R$ be continuous.  Then:
\begin{enumerate}
\item\label{itm:ch40-cmp:easy}
 $Q_L(h)\le\sup\{\tau\in\PSH_X(c_1(L)):\tau\le h\}$ on $X^{\An}$;
\item\label{itm:ch40-cmp:divisorial}
 every $\tau\in\PSH_X(c_1(L))$ with $\tau\le h$ satisfies
 $\tau\le Q_L(h)$ on $X^{\divis}$.  Consequently the raw supremum in
 \cref{itm:ch40-cmp:easy} coincides with $Q_L(h)$ on $X^{\divis}$.
\end{enumerate}
\end{proposition}

\begin{proof}
\cref{itm:ch40-cmp:easy}  Every
$\psi\in\cH^{\mathrm{gf}}_{\mathbf R}(L)$ lies in $\PSH_X(c_1(L))$, by
\Cref{def:bj-psh} applied to the constant net $\psi_i=\psi$; so every
competitor for $Q_L$ is a competitor for the psh envelope.

\cref{itm:ch40-cmp:divisorial}  Fix $\tau\in\PSH_X(c_1(L))$ with
$\tau\le h$, and fix $x\in X^{\divis}$.  If $\tau(x)=-\infty$ there is
nothing to prove, so assume $\tau(x)\in\mathbf R$.

\emph{Step 1: a Kodaira decomposition.}  Since $c_1(L)$ is big, there are
an ample $\mathbf Q$-class $\alpha$ and an effective
$\mathbf Q$-divisor $E$ on $X$ with
$c_1(L)=\alpha+c_1(E)$.  Indeed, let $A$ be a very ample divisor and
choose $0\ne\eta\in H^0(X,A)$ with zero divisor $Z$; the exact sequence
\[
 0\to\cO_X(mL-A)\to\cO_X(mL)\to\cO_Z(mL)\to0
\]
gives $h^0(mL-A)\ge h^0(mL)-h^0(\cO_Z(mL))$.  The first term grows like
$m^n$ because $L$ is big, the second at most like $m^{n-1}$ because
$Z$ has dimension $n-1$; hence $h^0(mL-A)>0$ for some $m$, and
writing $mL-A\sim E'\ge0$ gives
$c_1(L)=m^{-1}c_1(A)+m^{-1}c_1(E')$, which is the required decomposition.
Fix a nonzero section $s_E$ of a multiple of $E$ and write
$\log|s_E|$ for the corresponding generically finite Fubini--Study
function for $c_1(E)$, normalized so that $\log|s_E|\le0$.  Since
$x\in X^{\divis}$ has full support, $\log|s_E|(x)\in\mathbf R$.

\emph{Step 2: a Dini argument.}  By \Cref{def:bj-psh} there are
$\mathbf Q$-line bundles $L_i$ with $c_1(L_i)\to c_1(L)$ and a
decreasing net $\tau_i\in\cH^{\mathrm{gf}}_{\mathbf R}(L_i)$ with
$\tau_i\downarrow\tau$.  Fix $\varepsilon>0$.  The sets
$U_i:=\{\tau_i<h+\varepsilon\}$ are open, because $\tau_i$ is upper
semicontinuous and $h$ is continuous; they increase with $i$; and they
cover the compact space $X^{\An}$, since $\tau_i(y)\downarrow\tau(y)\le
h(y)<h(y)+\varepsilon$ for every $y$.  By compactness and directedness
there is $i_0$ with $\tau_i\le h+\varepsilon$ for all $i\ge i_0$.

\emph{Step 3: a competitor.}  Put $C:=\min_{X^{\An}}h$, which is finite
because $h$ is continuous on a compact space, and let
$\delta\in\mathbf Q\cap(0,1)$.  For $i\ge i_0$ set
\[
 \chi_{i,\delta}
 :=(1-\delta)\tau_i+\delta\bigl(\log|s_E|+C\bigr)-\varepsilon .
\]
This is a generically finite Fubini--Study function for the class
$(1-\delta)c_1(L_i)+\delta c_1(E)$, by
\Cref{prop:ch34-fs-operations}\,\cref{itm:ch34-ops:rational} and
\Cref{lem:gf-fs-operations}.  Moreover $\chi_{i,\delta}\le h$: indeed
$\log|s_E|\le0$ and $\tau_i\le h+\varepsilon$ give
\[
 \chi_{i,\delta}\le(1-\delta)(h+\varepsilon)+\delta C-\varepsilon
 = h-\delta\bigl(h+\varepsilon-C\bigr)\le h,
\]
because $h\ge C$.

\emph{Step 4: correcting the class.}  Let
$\gamma_{i,\delta}:=c_1(L)-(1-\delta)c_1(L_i)-\delta c_1(E)
=(1-\delta)\bigl(c_1(L)-c_1(L_i)\bigr)+\delta\alpha$.  Since $\alpha$ is
ample and $c_1(L_i)\to c_1(L)$, the class $\gamma_{i,\delta}$ is an
ample $\mathbf Q$-class for $i$ large.  Choose $N\ge1$ with
$N\gamma_{i,\delta}$ the class of a very ample line bundle and sections
$\sigma_1,\ldots,\sigma_r$ of it without common zero.  For every
$y\in X^{\An}$ some $\sigma_l$ is a unit in $\cO_{X,c(y)}$, so
$\log|\sigma_l|(y)=0$; hence $N^{-1}\max_l\log|\sigma_l|\equiv0$ and it
belongs to $\cH_{\mathbf Q}(\gamma_{i,\delta})$.  Adding it to
$\chi_{i,\delta}$ changes nothing and moves the class:
$\chi_{i,\delta}\in\cH^{\mathrm{gf}}_{\mathbf R}(L)$.  Therefore
$\chi_{i,\delta}\le Q_L(h)$.

\emph{Step 5: passing to the limit.}  Evaluating at $x$ and using
$\tau_i\ge\tau$,
\[
 Q_L(h)(x)\ \ge\ (1-\delta)\tau(x)
   +\delta\bigl(\log|s_E|(x)+C\bigr)-\varepsilon .
\]
All the quantities on the right are finite.  Letting $\delta\downarrow0$
and then $\varepsilon\downarrow0$ gives $Q_L(h)(x)\ge\tau(x)$.

The final assertion follows by combining \cref{itm:ch40-cmp:easy} and the
inequality just proved.
\end{proof}

\begin{remark}[The addendum's identity off the augmented base locus]
\label{rem:ch40-off-base-locus}
The addendum proves
\[
                       P_L(\phi_D)=Q_L(\phi_D)
\tag{40.8}\label{eq:bj-smooth-off-base}
\]
off the augmented base locus $\mathbf B_+(L)$; this is
\cite[Lemma~1.5]{BJ22b}.  The statement is genuinely stronger than
\Cref{prop:ch40-envelope-comparison}\,\cref{itm:ch40-cmp:divisorial}, because
it applies at all points of $X^{\An}$ whose centre avoids
$\mathbf B_+(L)$, and its proof uses the description of
$\mathbf B_+(L)$ as the intersection of the supports of the effective
parts of all Kodaira decompositions, \cite{ELMNP05}.  We do not
need it.  Every use made below is at a divisorial point, and every
divisorial point has full support, hence lies over the generic point of
$X$ and outside the analytification of any proper closed subset; so
\cref{itm:ch40-cmp:divisorial}, whose proof above is self-contained, is
exactly the required statement.  Recording this is not pedantry: the
augmented base locus is the one geometric object in this circle of ideas
whose behaviour is characteristic-sensitive in the literature, and the proof
below is free of it.
\end{remark}

\subsection{A rational obstacle on the projective plane, computed}
\label{sec:ch40-p2-example}

Everything above is explicit for the simplest obstacle that is not already
psh, and the resulting graded sequence turns out to be a single ideal raised
to powers.  We record the computation once and return to it three times
below: in \Cref{ex:ch40-p2-engines} to compute the two ideal engines on it,
in \Cref{ex:ch40-monomial-valuation} to check the valuation inequality on it
by hand, and in \Cref{ex:ch40-p2-envelope} to read off the envelope.

\begin{example}[The double point obstacle on $\mathbf P^2$]
\label{ex:ch40-p2-obstacle}
Let $X=\mathbf P^2$ with homogeneous coordinates $[X_0:X_1:X_2]$, let
$p=[0:0:1]$, and work in the affine chart $X_2\ne0$ with coordinates
$x=X_0/X_2$, $y=X_1/X_2$.  Let $I_p\subset\cO_X$ be the ideal of
$p$.  Take
\[
 L:=\cO_{\mathbf P^2}(1),
 \qquad
 S=\mathbf P^2\times\mathbf A^1,
 \qquad
 \mathfrak m:=I_p\cO_S+(\varpi)\subset\cO_S,
\]
so that $\mathfrak m$ is the ideal of the closed point $(p,0)\in S$ and,
in the chart, $\mathfrak m=(x,y,\varpi)$.  Let
$\mu:\cY:=\operatorname{Bl}_{(p,0)}S\to S$ be the blowup, with exceptional
divisor $\Xi\cong\mathbf P^2$, and set
\[
 D:=-2\Xi,\qquad \cL:=\mu^*L+D .
\]
Since $S$ is regular and $(p,0)$ is a closed point,
$\mathfrak m\cO_{\cY}=\cO_{\cY}(-\Xi)$ and
$\mu_*\cO_{\cY}(-k\Xi)=\mathfrak m^k$ for every $k\ge0$, the powers of
the maximal ideal of a regular local ring being integrally closed.

\emph{(i) The obstacle.}  Write $s:=v(I_p)\in[0,+\infty]$ for
$v\in X^{\An}$.  Then $\sigma(v)(\mathfrak m)=\min\{s,1\}$ by
\eqref{eq:ch40-gauss}, so
\[
 \phi_D=\log|\mathfrak m^2|=-2\min\{s,1\}
 =\max\{2\log|I_p|,\,-2\}.
\]
This function is \emph{not} $c_1(L)$-psh.  Indeed the central fibre of
$\cY$ is $\widetilde X\cup\Xi$, where
$\widetilde X\cong\operatorname{Bl}_p\mathbf P^2$ is the strict transform
of $X\times\{0\}$, and on $\widetilde X$ one has
$\theta_{\cY}|_{\widetilde X}=\pi^*H$ and
$D|_{\widetilde X}=-2E'$, where $H$ is the hyperplane class and $E'$
the exceptional curve.  Testing against the strict transform $C$ of a line
through $p$, whose class is $\pi^*H-E'$,
\[
 (\pi^*H-2E')\cdot(\pi^*H-E')=H^2+2(E')^2=1-2=-1<0,
\]
so $\theta_{\cY}+D$ is not relatively nef and
\Cref{lem:ch34-pl-positivity-well-posed} applies.  By contrast
$\phi_{-\Xi}=\max\{\log|I_p|,-1\}$ \emph{is} $c_1(L)$-psh: on
$\widetilde X$ the restriction of $\theta_{\cY}-\Xi$ is the base point
free class $\pi^*H-E'$, and on $\Xi\cong\mathbf P^2$ it is
$-\Xi|_{\Xi}=\cO_{\mathbf P^2}(1)$.

\emph{(ii) The relative sections.}  By the projection formula,
\[
 \mu_*\cO_{\cY}(m\cL)
 =\cO_S(m)\otimes\mathfrak m^{2m},
 \qquad
 M_m=H^0\bigl(S,\cO_S(m)\otimes\mathfrak m^{2m}\bigr),
\]
where $\cO_S(m):=\operatorname{pr}_X^*\cO_{\mathbf P^2}(m)$.  A section is
a sum $\sum_kF_k\varpi^k$ with $F_k$ a form of degree $m$; in the
chart, $F_k/X_2^m=f_k(x,y)$ has degree $\le m$, and
$\mathfrak m^{2m}=\bigoplus_k(x,y)^{2m-k}\varpi^k$ is a monomial ideal, so
the condition is $f_k\in(x,y)^{2m-k}$ for every $k$.  As
$\deg f_k\le m$, this forces $k\ge m$, and for a monomial of
$x,y$-degree $l\le m$ the least admissible $k$ is $2m-l$.  Hence
$M_m$ is free over $k[\varpi]$ of rank $\binom{m+2}{2}=h^0(\cO(m))$,
with basis
\[
 \bigl\{\varpi^{2m-l}X_0^iX_1^jX_2^{m-l}\ :\
 i+j=l,\ 0\le l\le m\bigr\}.
\]

\emph{(iii) The ambient ideal.}  Twisting down by $mL$ and taking the
image as in \eqref{eq:bj-smooth-ambient-ideal}, the ideal $J_m$ is
generated by the elements $\varpi^{2m-l}\cdot\{\text{degree-}m\text{ part of
}I_p^{\,l}\}$.  Since $I_p^{\,l}$ is generated in degree $l\le m$, its
degree-$m$ part generates the sheaf $I_p^{\,l}$, whence
\[
 J_m=\sum_{l=0}^{m}\varpi^{2m-l}\,I_p^{\,l}
    =\varpi^m\sum_{l=0}^{m}\varpi^{m-l}I_p^{\,l}
    =\varpi^m\mathfrak m^m
    =(\varpi\mathfrak m)^m .
\]
Thus the graded sequence is the sequence of powers of the single ideal
$J_1=\varpi\mathfrak m$, and in particular
$J_mJ_r=J_{m+r}$, so \eqref{eq:bj-smooth-graded-ideals} holds with
equality.

\emph{(iv) The two gradings.}  In the weight decomposition of
\Cref{lem:ch40-equivariant-ideal-fs}\,\cref{itm:ch40-eqfs:decomp},
\[
 \mathfrak b_{m,k}=
 \begin{cases}
 0,&k<m,\\
 I_p^{\,2m-k},&m\le k\le 2m,\\
 \cO_X,&k\ge2m,
 \end{cases}
\]
so the weight decomposition of $J_m$ has exactly $m+1$ distinct nonzero
terms: the second grading is not trivial, and it is not bounded
independently of the degree $m$.  Each $\mathfrak b_{m,k}(mL)
=I_p^{\,2m-k}(m)$ with $2m-k\le m$ is globally generated, confirming
\Cref{lem:ch40-J-basic}\,\cref{itm:ch40-J:coherent} by inspection.  The
resulting Fubini--Study presentation is
\[
 q_m=\frac1m\max_{0\le l\le m}
     \bigl\{l\log|I_p|-(2m-l)\bigr\}
    =-2+\frac1m\max_{0\le l\le m}l\,(1-s),
\]
a weighted section system of degree $m$ with the $m+1$ integer weights
$-(2m-l)$.

\emph{(v) The base-ideal limit.}  Evaluating the maximum: for $s\le1$ it
is attained at $l=m$ and for $s\ge1$ at $l=0$.  Hence
\[
 q_m=-\bigl(1+\min\{s,1\}\bigr)
 \qquad\text{for every }m\ge1,
\]
so the increasing net of \Cref{prop:ch40-Q-limit} is constant here, and
\[
 Q=Q_L(\phi_D)=-\bigl(1+\min\{s,1\}\bigr)
   =\max\{\log|I_p|,-1\}-1
   =\phi_{-\Xi}-1 .
\]
One checks $Q\le\phi_D$ directly:
$-(1+\min\{s,1\})\le-2\min\{s,1\}$ is the inequality
$\min\{s,1\}\le1$, with equality exactly when $s\ge1$.

\emph{(vi) The relative base ideal.}  From
$\varpi\cO_{\cY}=\cO_{\cY}(-\widetilde X-\Xi)$ and
$\mathfrak m\cO_{\cY}=\cO_{\cY}(-\Xi)$ we get
$J_m\cO_{\cY}=\cO_{\cY}(-m\widetilde X-2m\Xi)$, and comparison with
\eqref{eq:bj-smooth-base-ideal-equality} gives
\[
 \mathfrak a_m=\cO_{\cY}(-m\widetilde X):
\]
the relative base locus of $|m\cL|$ is the strict transform of the central
fibre of $S$, with multiplicity $m$.  Consistently,
$m^{-1}\phi_{\mathfrak a_m}=-\bigl(1-\min\{s,1\}\bigr)$ and
\eqref{eq:bj-smooth-obstacle-identity} reads
\[
 -\bigl(1+\min\{s,1\}\bigr)
 =-2\min\{s,1\}-\bigl(1-\min\{s,1\}\bigr),
\]
which is an identity.  In particular $Q(v_{\triv})=-1<0=\phi_D(v_{\triv})$:
the envelope is strictly below the obstacle already at the trivial
valuation, because a psh function in the class $\cO(1)$ cannot afford a
pole of order two at $p$ without paying a constant.
\end{example}

\begin{remark}[What the example already shows]
\label{rem:ch40-p2-lessons}
Three features of \Cref{ex:ch40-p2-obstacle} are typical and one is not.

Typical: the obstacle is a PL function which is \emph{not} psh, the
Fubini--Study envelope $Q$ is strictly below it on a nonempty open set,
and the relative base ideal $\mathfrak a_m$ is a nontrivial divisorial
ideal supported on the central fibre.  These are the three phenomena the
machinery exists to handle.

Also typical: the second grading is genuinely present, with $m+1$ distinct
weight pieces in degree $m$.  Any argument that treats $J_m$ as if it
were extended from $X$, that is, as if
\Cref{prop:ch40-two-gradings-here}\,\cref{itm:ch40-grad:homogeneous} applied,
computes with $\mathfrak b_{m,m}=I_p^{\,m}$ alone and obtains
$-2+(1-s)$ rather than the correct $-1-\min\{s,1\}$; the two agree only
for $s\le1$.

Not typical: the graded sequence is the sequence of powers of a single
ideal, so the increasing net $(q_m)$ is constant.  This happens because
the obstacle is attached to a divisor $D$ that is $\mu$-semiample after
subtracting the base locus; for a general PL obstacle the net is strictly
increasing and the limit $Q$ is not attained at any finite stage.
\Cref{ex:ch40-drift-real} exhibits the general behaviour.
\end{remark}

\section{The two ideal engines}
\label{sec:ch40-engines}

The graded sequence $J_\bullet$ produced in \Cref{sec:ch40-base-ideals}
has one defect and one virtue.  The virtue is that its normalized functions
$q_m$ converge to $Q$, which by
\Cref{prop:ch40-envelope-comparison}\,\cref{itm:ch40-cmp:divisorial} is the
right limit.  The defect is that the convergence is \emph{increasing}, and
\Cref{def:bj-psh} recognizes only decreasing limits.  The role of an ideal
engine is to repair the direction of the convergence without moving the
limit.

Formally, we ask for an assignment $m\mapsto B_m$ of coherent ideals on
$S$ with four properties: $B_m\supseteq J_m$, so that $u_m\ge q_m$ and
the limit is not too small; $B_{m\ell}\subseteq B_m^{\,\ell}$, so that
$u_m$ decreases; a relative generation statement with a twist independent
of $m$, so that $u_m$ is a Fubini--Study function in a class converging
to $c_1(L)$; and a valuation estimate with an error $o(m)$, so that the
limit is not too large.  The first, second and fourth are treated in this
section; the third, which is the longest, in \Cref{sec:ch40-stages}.

\begin{definition}[The stage ideals]
\label{def:ch40-stage-ideals}
For $m\ge1$ put
\[
 B_m:=
 \begin{cases}
 \mathcal J_S(J_\bullet^{\,m}),
      &\operatorname{char}k=0,\\[2pt]
 \tau_S(J_\bullet^{\,m}),
      &\operatorname{char}k=p>0,
 \end{cases}
 \qquad
 u_m:=\frac1m\phi_{B_m},
\]
where $\mathcal J_S$ is the asymptotic multiplier ideal
\Cref{def:appC-asymptotic-multiplier} and $\tau_S$ is the asymptotic test
ideal \Cref{def:appC-asymptotic-test}.  Both are defined because $S$ is
smooth over $k$ --- hence, in characteristic $p$, regular and
$F$-finite --- and because $J_\bullet$ is a nonzero graded sequence by
\Cref{lem:ch40-J-basic}.
\end{definition}

\subsection{What both engines supply}
\label{sec:ch40-engine-package}

\begin{proposition}
\label{prop:ch40-engine-package}
In either characteristic, and for all $m,\ell\ge1$:
\begin{enumerate}
\item\label{itm:ch40-eng:contain}
 $J_m\subseteq B_m\subseteq\cO_S$, and $B_m\ne0$;
\item\label{itm:ch40-eng:subadd}
 $B_{m\ell}\subseteq B_m^{\,\ell}$;
\item\label{itm:ch40-eng:equivariant}
 $B_m$ is $\mathbf G_m$-equivariant;
\item\label{itm:ch40-eng:sandwich}
 $u_{m\ell}\le u_m$ and
 \[
  Q\le u_m\le0\qquad\text{on }X^{\An};
 \]
\item\label{itm:ch40-eng:valuation}
 for every divisorial valuation $w$ centred on $S$,
 \[
  w(B_m)\ \ge\ m\,w(J_\bullet)-A_S(w),
  \qquad
  w(J_\bullet):=\inf_{q\ge1}\frac{w(J_q)}{q}.
 \]
\end{enumerate}
\end{proposition}

\begin{proof}
\cref{itm:ch40-eng:contain}  In characteristic zero this is
\Cref{prop:appC-asymptotic-multiplier-properties}\,(1); in characteristic
$p$ it is the first inclusion of
\eqref{eq:appC-test-subadditivity}.  Both ideals are contained in
$\cO_S$ by construction, and both contain $J_m\ne0$.

\cref{itm:ch40-eng:subadd}  In characteristic zero this is
\eqref{eq:appC-asymptotic-subadditivity}, which rests on the
Demailly--Ein--Lazarsfeld subadditivity theorem
\Cref{thm:appC-multiplier-subadditivity}; in characteristic $p$ it is the
second inclusion of \eqref{eq:appC-test-subadditivity}, from
\cite{MustataTest}.

\cref{itm:ch40-eng:equivariant}  For $\lambda\in k^\times$ let
$g_\lambda$ be the automorphism of $S=X\times\mathbf A^1$ given by
$\varpi\mapsto\lambda\varpi$.  Each $J_q$ is $g_\lambda$-stable by
\Cref{lem:ch40-J-basic}\,\cref{itm:ch40-J:coherent}, and both constructions
of \Cref{def:ch40-stage-ideals} commute with isomorphisms of regular
($F$-finite) schemes: the multiplier ideal because
\eqref{eq:appC-multiplier-resolution} may be computed on the transported log
resolution, the test ideal because \eqref{eq:appC-test-trace} is built from
the Grothendieck trace, which is natural.  Hence
$g_\lambda^*B_m=B_m$ for every $\lambda\in k^\times$.  Now write, on an
affine chart, $B_m\ni f=\sum_{k}f_k\varpi^k$ with $f_k$ independent of
$\varpi$.  Then $\sum_k\lambda^kf_k\varpi^k\in B_m$ for every
$\lambda\in k^\times$.  Choosing $N+1$ distinct values of $\lambda$,
where $N$ bounds the $\varpi$-degree of $f$, and inverting the
Vandermonde matrix --- possible since $k$ is infinite, being algebraically
closed --- gives $f_k\varpi^k\in B_m$ for every $k$.  Thus $B_m$ is a
graded ideal, that is, $\mathbf G_m$-equivariant.

\cref{itm:ch40-eng:sandwich}  From \cref{itm:ch40-eng:subadd} and the order
reversal of \Cref{warn:ch40-signs},
\[
 \phi_{B_{m\ell}}\le\phi_{B_m^{\,\ell}}=\ell\,\phi_{B_m},
\]
using \eqref{eq:ch34-log-ideal-semiring}; dividing by $m\ell$ gives
$u_{m\ell}\le u_m$.  The upper bound $u_m\le0$ is
$B_m\subseteq\cO_S$.  For the lower bound, fix $m$; for every
$\ell\ge1$, \cref{itm:ch40-eng:contain} in degree $m\ell$ gives
$q_{m\ell}\le u_{m\ell}\le u_m$, and letting $\ell\to\infty$ along
divisibility, \Cref{prop:ch40-Q-limit}\,\cref{itm:ch40-Q:sup} gives
$Q\le u_m$.  No finiteness of any of these values is assumed; all
inequalities are inequalities in $[-\infty,0]$.

\cref{itm:ch40-eng:valuation}  In characteristic zero this is
\eqref{eq:appC-asymptotic-discrepancy}.  In characteristic $p$ it follows
from \Cref{lem:bj-smooth-canton} below, applied with $I=J_q$ and
$t=m/q$ for $q$ sufficiently divisible: asymptotic stabilization
\eqref{eq:appC-test-stabilization} gives $B_m=\tau_S(J_q^{m/q})$, and
\[
 w(B_m)\ \ge\ \frac mq\,w(J_q)-A_S(w)\ \ge\ m\,w(J_\bullet)-A_S(w),
\]
the last step because $q^{-1}w(J_q)\ge w(J_\bullet)$ by definition of the
infimum.
\end{proof}

\begin{remark}[Two theorems with the same statement and no common step]
\label{rem:ch40-subadditivity-two-proofs}
Item \cref{itm:ch40-eng:subadd} is the place where the two arguments diverge
most sharply, and the divergence is invisible in the statement.

In characteristic zero, \Cref{thm:appC-multiplier-subadditivity} is proved by
forming the external product of $\mathfrak a$ and $\mathfrak b$ on
$S\times S$, where the multiplier ideal of a product of pullbacks is the
product of the pullbacks of the multiplier ideals, and then restricting to
the diagonal.  The restriction theorem for multiplier ideals is a corollary
of Nadel vanishing, hence of Kodaira vanishing, hence unavailable in
characteristic $p$; and the smoothness of the diagonal of $S\times S$
--- that is, the smoothness of $S$ --- is used, which is why the statement
fails on singular varieties in any characteristic.  See \cite{DEL00}.

In characteristic $p$ there is no product-and-restrict argument.  The
inclusion $\tau(\mathfrak a_\bullet^{m\ell})\subseteq
\tau(\mathfrak a_\bullet^{m})^{\ell}$ is proved directly from the trace
description \eqref{eq:appC-test-trace}: a Frobenius root of an element of
$\mathfrak a_q^{\lceil m\ell p^e/q\rceil}$ is manipulated into a product of
$\ell$ Frobenius roots of elements of
$\mathfrak a_q^{\lceil mp^e/q\rceil}$, the rounding being absorbed because
$\lceil m\ell t\rceil\ge\ell\lceil mt\rceil-(\ell-1)$ and the resulting
finite discrepancy is killed by the asymptotic stabilization.  See
\cite{MustataTest} and, for the underlying finite-characteristic
formalism, \cite{ST12}.

The two proofs share the conclusion and nothing else.  A reader who
remembers only ``subadditivity holds'' will not be misled about this
chapter; a reader who remembers ``subadditivity holds by restriction to the
diagonal'' will be misled about \Cref{chap:bj-alterations}, where the
absence of a diagonal argument on singular sources is one of the reasons the
smooth case has to be settled first.
\end{remark}

\subsection{Characteristic zero}
\label{sec:ch40-char-zero}

In characteristic zero the construction of $B_m$ rests on the existence of
log resolutions of the pairs $(S,J_q)$, which is
\Cref{thm:appB-resolution}, and on local vanishing.  Both are recorded in
\Cref{sec:appC-algebraic-multiplier} together with the valuative description
\Cref{thm:appC-valuative-multiplier}, from which the containment
$J_m\subseteq B_m$ is immediate: for $f\in J_m$ and $v$ divisorial,
$v(f)\ge v(J_m)\ge(m/q)v(J_q)$ by gradedness, and
$A_S(v)>0$ because $S$ is regular, so the strict inequality
$v(f)+A_S(v)>(m/q)v(J_q)$ required by
\eqref{eq:appC-valuative-multiplier} holds.  It is worth pausing on this
last point: the containment $\mathfrak a_m\subseteq\mathcal J(\mathfrak
a_\bullet^m)$ is not a formal consequence of monotonicity, and it uses the
positivity of the log discrepancy on a regular scheme, that is, exactly the
smoothness hypothesis.

\subsection{Positive characteristic}
\label{sec:ch40-char-p}

Assume now that $\operatorname{char}k=p>0$.  The regular scheme $S$ is
$F$-finite.  For a nonzero graded ideal sequence $J_\bullet$, its
asymptotic test ideals satisfy
\[
 J_m\subset\tau_S(J_\bullet^m),\qquad
 \tau_S(J_\bullet^{m\ell})
       \subset\tau_S(J_\bullet^m)^\ell.
\tag{40.9}\label{eq:bj-test-ideal-properties}
\]
The second inclusion is subadditivity.

Define
\[
 B_m:=\tau_S(J_\bullet^m),\qquad
 u_m:=\frac1m\phi_{B_m}.
\tag{40.10}\label{eq:bj-smooth-test-stages}
\]
Ideal inclusion reverses valuation order after negation.  Hence $u_m$
decreases under divisibility.  For every multiple $m\ell$,
\[
 \frac1{m\ell}\phi_{J_{m\ell}}
       \le u_{m\ell}\le u_m .
\]
Taking the cofinal limit on the left, and using $B_m\subset\cO_S$, gives
\[
                            Q\le u_m\le0
                            \quad\text{on }X^{\An}.
\tag{40.11}\label{eq:bj-smooth-test-sandwich}
\]
No finiteness of these values is assumed.

The displayed assertions are
\Cref{prop:ch40-engine-package}\,\cref{itm:ch40-eng:contain},
\cref{itm:ch40-eng:subadd} and \cref{itm:ch40-eng:sandwich}, and
\eqref{eq:bj-smooth-test-stages} is the positive-characteristic instance of
\Cref{def:ch40-stage-ideals}.  Nothing in this subsection assumes that
$J_\bullet$ comes from an obstacle; the sole hypotheses are that $S$ is
regular and $F$-finite and that $J_\bullet$ is a nonzero graded
sequence.

\begin{example}[The powers of a smooth parameter]
\label{ex:ch40-smooth-parameter}
Let $R=k[x]$, where $k$ is an $F$-finite field of characteristic
$p>0$, and consider the graded sequence $J_q=(x^q)$ on
$\Spec R$.  For every $m\geq1$, stabilization may be computed using any
sufficiently divisible $q$, and the standard test-ideal formula for a
smooth principal divisor \cite{ST12} gives
\[
       \tau_R\bigl(J_q^{m/q}\bigr)
          =\tau_R\bigl((x)^m\bigr)=(x^m).
\]
Consequently
\[
       B_m=(x^m),\qquad
       u_m=\frac1m\log|(x^m)|=\log|(x)|.
\]
The complete base-ideal limit is the same function, since
$q^{-1}\log|J_q|=\log|(x)|$.  Thus the \emph{left} inequality in
\eqref{eq:bj-smooth-test-sandwich} is an equality in this
simple-normal-crossings model: $u_m=Q$ for every $m$.  The right
inequality $u_m\le0$ is of course strict wherever $x$ is not a unit ---
at $w=\ord_0$ one has $u_m(w)=-w(x)=-1<0$.  In particular, the normalization by $m^{-1}$ introduces
no hidden Frobenius factor.
\end{example}

\begin{remark}[Why the last sentence of \Cref{ex:ch40-smooth-parameter} is
worth checking]
\label{rem:ch40-no-frobenius-factor}
The trace description \eqref{eq:appC-test-trace} involves the exponent
$\lceil tp^e\rceil$, and it is a natural worry that the resulting ideal
carries a factor $p^e$ which fails to cancel after dividing by $m$.  It
does cancel, and \Cref{ex:ch40-smooth-parameter} is the unit test: the
answer $(x^m)$ is exactly the $m$-th power, with no $p$-adic
distortion, so $u_m=\log|(x)|$ is independent of $m$ and of $p$.  The
same test performed with the convention $p^e-1$ in place of $p^e$, which
occurs in parts of the literature on Cartier algebras, produces
$(x^{m})$ as well but only after the rounding has been tracked; this is
the point of the convention remark following
\Cref{def:appC-test-ideal}, and it is why that convention is fixed once and
for all in \Cref{app:ideals} rather than chosen locally.
\end{remark}

\subsection{The engines on the plane example}
\label{sec:ch40-p2-engines}

\begin{example}[The stage ideals of \Cref{ex:ch40-p2-obstacle}]
\label{ex:ch40-p2-engines}
Keep the notation of \Cref{ex:ch40-p2-obstacle}, so that
$J_m=(\varpi\mathfrak m)^m$ on $S=\mathbf P^2\times\mathbf A^1$ with
$\mathfrak m$ the ideal of $(p,0)$.  We compute $B_m$ in both
characteristics.  Since the two constructions are local, it is enough to
work at the three kinds of point of $S$.

\emph{Points with $\varpi\ne0$.}  There $\mathfrak m=\cO_S$ and
$\varpi$ is a unit, so $J_m=\cO_S$ and $B_m=\cO_S$.

\emph{Points $(q,0)$ with $q\ne p$.}  There $\mathfrak m=\cO_S$, so
$J_m=(\varpi^m)$ is principal, generated by an element whose divisor is
smooth.  In a regular system of parameters containing $\varpi$ this ideal
is monomial, and both engines return
$B_m=(\varpi^m)$: in characteristic zero by the simple normal crossings
formula \eqref{eq:appC-snc-formula} with $\lfloor m\rfloor=m$, and in
characteristic $p$ by \Cref{ex:ch40-smooth-parameter} applied in the
variable $\varpi$.

\emph{The point $(p,0)$.}  Here $(x,y,\varpi)$ is a regular system of
parameters and
\[
 J_m=\varpi^m(x,y,\varpi)^m
\]
is a monomial ideal.  Its Newton polyhedron is
\[
 P=\bigl\{(u_1,u_2,u_3)\in\mathbf R^3_{\ge0}\ :\
 u_3\ge1,\ u_1+u_2+u_3\ge2\bigr\},
\]
the translate by $(0,0,1)$ of the Newton polyhedron of
$\mathfrak m$.  By Howald's formula, valid for multiplier ideals of
monomial ideals in characteristic zero and for generalized test ideals of
monomial ideals in characteristic $p$ --- the latter is the monomial case
of \cite{HY03}, surveyed in \cite{ST12} --- the ideal
$B_m$ is generated by the monomials $x^{u_1}y^{u_2}\varpi^{u_3}$ with
$(u_1+1,u_2+1,u_3+1)$ in the topological interior of $mP$, that is, with
\[
 u_3+1>m
 \quad\text{and}\quad
 (u_1+1)+(u_2+1)+(u_3+1)>2m,
\]
the constraints $u_i+1>0$ being automatic.  For integral $u$ these read
$u_3\ge m$ and $u_1+u_2+u_3\ge2m-2$, whence
\[
 B_m=\varpi^m\,\mathfrak m^{\,m-2}
 \qquad(m\ge2),
 \qquad B_1=\varpi\,\cO_S .
\]
The same formula is obtained by factoring out the principal part:
$B_m=\varpi^m\,\mathcal J_S(\mathfrak m^m)$ with
$\mathcal J_S(\mathfrak m^m)=\mathfrak m^{m-2}$, the exponent being
$m-\dim S+1=m-2$.

\emph{The stage functions.}  Let $m\ge2$, which is no restriction since
the divisibility net may be indexed cofinally by such $m$.  With
$s=v(I_p)$ as before, $\sigma(v)(B_m)=m+(m-2)\min\{s,1\}$, so
\[
 u_m=-\Bigl(1+\Bigl(1-\frac2m\Bigr)\min\{s,1\}\Bigr),
 \qquad
 u_m-Q=\frac2m\min\{s,1\}\ \ge0 .
\]
The net $(u_m)$ decreases to $Q$, as
\Cref{prop:ch40-engine-package}\,\cref{itm:ch40-eng:sandwich} predicts, and
the gap is exactly $2/m$ at every point with $s\ge1$.  It is strictly
positive for every finite $m$: no stage of the construction computes the
envelope, and the limit is genuinely needed.

\emph{Relative generation without a twist.}  The weight pieces of $B_m$
are $I_p^{\,2m-2-k}$ for $m\le k\le2m-2$, and
$I_p^{\,j}(m)$ is globally generated for $j\le m-2$; so
$B_m(mL)$ is already generated relative to $\mathbf A^1$, and the twist
$A$ of \Cref{sec:ch40-stages} may be taken trivial here.  This is not an
accident of the example: on $X=\mathbf P^n$ with $H=\cO(1)$ and
$d=n+1$ one has $A=K_X+(n+1)H=\cO_X$, so the class drift vanishes
identically.  \Cref{ex:ch40-drift-real} shows that this is a feature of
projective space and not of the construction.
\end{example}

\subsection{Where the two engines genuinely differ}
\label{sec:ch40-engines-differ}

The plane example produced the same ideal in both characteristics, because
the ideals involved were monomial and Howald's description is
characteristic-free.  That agreement is special.  The next example, which is
the smallest classical one, exhibits a graded sequence on a smooth surface
for which the asymptotic multiplier ideal and the asymptotic test ideal are
different sheaves.

\begin{example}[The cusp: the two engines separate]
\label{ex:ch40-cusp}
Let $k$ be algebraically closed, let $V:=\mathbf A^2_k=\Spec k[x,y]$, and
let
\[
 f:=x^2+y^3,
\]
the local equation of a cuspidal curve at the origin $o$.  Consider the
graded sequence
\[
 \mathfrak c_q:=\bigl(f^{\lceil 4q/5\rceil}\bigr),
 \qquad q\ge1,
\]
which is graded because
$\lceil a\rceil+\lceil b\rceil\ge\lceil a+b\rceil$.  For $q\in5\mathbf Z$
one has $\lceil4q/5\rceil=4q/5$ and hence
$\mathfrak c_q^{\,m/q}=(f^{4m/5})$ as a formal power; since
$5\mathbf Z$ is cofinal for divisibility, the stabilization
\eqref{eq:appC-asymptotic-stabilization}, respectively
\eqref{eq:appC-test-stabilization}, gives
\[
 \mathcal J_V(\mathfrak c_\bullet^{\,m})=\mathcal J_V\bigl(f^{4m/5}\bigr),
 \qquad
 \tau_V(\mathfrak c_\bullet^{\,m})=\tau_V\bigl(f^{4m/5}\bigr).
\]
We compare the two at $m=1$.

\emph{Characteristic zero: the log canonical threshold is $5/6$.}  Resolve
the cusp by three point blowups.  Blowing up $o$ produces $E_1$ with
$\ord_{E_1}(f)=2$ and $A_V(\ord_{E_1})=2$; in the chart
$x=st,\ y=t$ one has $f=t^2(s^2+t)$, so the strict transform is smooth
and tangent to $E_1$.  Blowing up the tangency point produces $E_2$ with
$\ord_{E_2}(f)=1+2=3$ and $A_V(\ord_{E_2})=2+1=3$.  On the resulting
surface the strict transform, $E_1$ and $E_2$ pass through a common
point; blowing it up produces $E_3$ with
\[
 \ord_{E_3}(f)=1+2+3=6,
 \qquad
 A_V(\ord_{E_3})=2+\ord_{E_3}(E_1'+2E_2)=2+3=5,
\]
and the total transform is now simple normal crossings.  By
\Cref{thm:appC-valuative-multiplier}, $\mathcal J_V(f^{c})=\cO_V$ if and
only if $c\,\ord_F(f)<A_V(\ord_F)$ for every prime divisor $F$ over
$V$, and on a log resolution it suffices to test the finitely many
components of the total transform.  The ratios are
$2/2$, $3/3$, $5/6$ and, for the strict transform of the curve,
$1/1$; hence
\[
 \lct_o(f)=\tfrac56,
 \qquad\text{and in particular}\qquad
 \mathcal J_V\bigl(f^{4/5}\bigr)=\cO_V,
\]
because $4/5<5/6$.

\emph{Characteristic five: the $F$-pure threshold is $4/5$.}  Let now
$\operatorname{char}k=5$.  The $F$-pure threshold of the cusp is
\[
 \operatorname{fpt}_o(x^2+y^3)=\frac56-\frac1{6p}
 \ \ (p\equiv2\bmod 3),
 \qquad
 \operatorname{fpt}_o(x^2+y^3)=\frac56
 \ \ (p\equiv1\bmod3),
\]
for $p\ge5$; this is the explicit computation
\cite[Example~4.3]{MTW05}, and it is the one external input of this example.
For $p=5$, which satisfies $5\equiv2\bmod3$, the first value is
$5/6-1/30=4/5$.  By definition of the $F$-pure threshold,
$\tau_V(f^{c})=\cO_V$ for $c<4/5$, and by right continuity of
$c\mapsto\tau_V(f^c)$, \cite[Theorem~6.6]{ST12}, one has
$\tau_V(f^{4/5})=\tau_V(f^{4/5+\varepsilon})\ne\cO_V$ for small
$\varepsilon>0$.  Thus
\[
 \tau_V\bigl(f^{4/5}\bigr)\subsetneq\cO_V .
\]

\emph{Conclusion.}  For the same graded sequence $\mathfrak c_\bullet$ and
the same exponent $m=1$, the characteristic-zero engine returns the unit
ideal and the characteristic-five engine returns a proper ideal, contained
in the maximal ideal at $o$.  Their ideal functions already differ at the
divisorial valuation of the exceptional curve of the blowup of $o$: the
multiplier-ideal value is $0$, whereas the test-ideal value is strictly
negative.  In particular no identity of the form
``$\tau=\mathcal J$'' is available, and none is used.
\end{example}

\begin{warning}{Reduction modulo \texorpdfstring{$p$}{p} is not a bridge}
\label{warn:ch40-no-reduction}
It is tempting to argue as follows: the characteristic-zero proof works;
test ideals agree with multiplier ideals after reduction modulo $p$ by
\Cref{thm:appC-hara-yoshida}; therefore the characteristic-zero proof works
in characteristic $p$.  Every clause of that sentence is false as an
argument, for three separate reasons, and \Cref{ex:ch40-cusp} illustrates
all three.

First, \Cref{thm:appC-hara-yoshida} concerns \emph{one fixed} datum
$(V,\mathfrak a,c)$ spread over a finitely generated
$\mathbf Z$-algebra, and asserts equality only on a suitable dense open
set of closed fibres.  The cusp example shows exactly this excluded-fibre
caveat: at the fixed exponent $c=4/5$, the characteristic-five test ideal
is proper whereas the characteristic-zero multiplier ideal is trivial.  It
does \emph{not} exclude every prime $p\equiv2\bmod3$: for such $p>5$,
\[
 \operatorname{fpt}_o(x^2+y^3)=\frac56-\frac1{6p}>\frac45,
\]
so this particular exponent no longer separates the two ideals.

Second, the datum here is not fixed: the ideals $J_q$ depend on $q$, the
exponent depends on $m$, and the conclusion must hold for all $m$
simultaneously.  Even for a single graded sequence there is no ``uniform in
$m$'' form of \Cref{thm:appC-hara-yoshida}, and the envelope argument
consumes infinitely many $m$.

Third, and most simply, the ground field of this chapter is a fixed
algebraically closed field of characteristic $p$, typically not obtained by
reduction of anything.  There is no characteristic-zero model to reduce.

The correct statement of the relationship is the one recorded in
\Cref{check:appC-characteristic-firewall}: two engines, run separately, with
the same four output properties.
\end{warning}

\section{The fixed twist and the exact class error}
\label{sec:ch40-stages}

Put $n=\dim X$.  Choose a very ample globally generated line bundle
$H$, an integer $d\ge n+1$, and define the fixed twist
\[
                              A:=K_X+dH.
\tag{40.12}\label{eq:bj-smooth-fixed-twist}
\]
We write $H$ and $A$ also for their pullbacks
$\operatorname{pr}_X^*H$ and $\operatorname{pr}_X^*A$ to $S$; the
morphism $\rho:S\to\mathbf A^1$ is projective with $n$-dimensional
fibres, and $H$ is $\rho$-very ample and $\rho$-globally generated.

The object of this section is that $B_m(mL+A)$ is generated relative to
$\mathbf A^1$, with \emph{the same} $A$ for every $m$.  Three features
of \eqref{eq:bj-smooth-fixed-twist} deserve to be noted before the proofs.
The bundle $A$ does not depend on $m$, on $q$, on the obstacle, or on
the characteristic.  It does depend on $X$ and on the auxiliary choices
$H$ and $d$, and on nothing else; in particular it does not depend on
$L$, which is why the same $A$ serves for every obstacle in every big
class.  And the canonical bundle appears because both vanishing arguments
below produce a sheaf of the shape $\cO(K+\text{positive})$, in
characteristic zero through Kawamata--Viehweg vanishing and in
characteristic $p$ through the Grothendieck trace; the identity
\[
 K_S=\operatorname{pr}_X^*K_X,
 \qquad
 K_{\mathbf A^1}\simeq\cO_{\mathbf A^1},
\]
which is \eqref{eq:bj-smooth-canonical-product} below, is what allows the
twist to be chosen on $X$ rather than on $S$.

\subsection{Characteristic zero: relative generation by vanishing}
\label{sec:ch40-generation-char-zero}

\begin{proposition}
\label{prop:ch40-multiplier-uniform-generation}
Assume $\operatorname{char}k=0$ and let $A$ be as in
\eqref{eq:bj-smooth-fixed-twist}.  Then for every $m>0$ the sheaf
$B_m(mL+A)$ is generated relative to $\mathbf A^1$, and consequently
\[
 u_m\in\cH^{\mathrm{gf}}_{\mathbf Q}\bigl(L+A/m\bigr).
\]
The line bundle $A$ is independent of $m$.
\end{proposition}

\begin{proof}
Fix $m$ and choose $q$, divisible by $m$ and sufficiently divisible,
so that $B_m=\mathcal J_S(J_q^{\,m/q})$ by
\eqref{eq:appC-asymptotic-stabilization}.  Let
\[
 \nu:T\longrightarrow S,
 \qquad
 J_q\cO_T=\cO_T(-F),
\]
be a log resolution of the pair $(S,J_q)$, which exists by
\Cref{thm:appB-resolution} since $\operatorname{char}k=0$; thus $T$ is
smooth, $\nu$ is projective birational, and
$F+\operatorname{Exc}(\nu)$ has simple normal crossings support.  By
\eqref{eq:appC-multiplier-resolution},
\[
 B_m=\nu_*\cO_T\Bigl(K_{T/S}-\bigl\lfloor\tfrac mqF\bigr\rfloor\Bigr).
\]

\emph{Step 1: the higher direct images.}  Put
$G:=B_m(mL+A)$ and fix $1\le i\le n$.  Local vanishing,
\cite[Theorem~9.4.1]{PAG2}, gives
$R^j\nu_*\cO_T(K_{T/S}-\lfloor\tfrac mqF\rfloor)=0$ for $j>0$; by the
projection formula and the degeneration of the Leray spectral sequence,
\[
 R^i\rho_*\bigl(G(-iH)\bigr)
 =R^i(\rho\nu)_*\,\cO_T\Bigl(
   K_{T/S}-\bigl\lfloor\tfrac mqF\bigr\rfloor
   +\nu^*(mL+A-iH)\Bigr).
\]

\emph{Step 2: rewriting the twist.}  Since $K_S=\operatorname{pr}_X^*K_X$
and $A=K_X+dH$ on $X$, we have $\nu^*A=\nu^*K_S+d\,\nu^*H$ and
therefore
\[
 K_{T/S}+\nu^*A=K_T+d\,\nu^*H .
\]
Writing $\{\cdot\}$ for the fractional part, the divisor inside the direct
image becomes
\[
 K_T+P+\Delta,
 \qquad
 P:=\frac mq\bigl(\nu^*(qL)-F\bigr)+(d-i)\nu^*H,
 \qquad
 \Delta:=\Bigl\{\tfrac mqF\Bigr\}.
\]
Indeed $-\lfloor\tfrac mqF\rfloor+m\nu^*L
=\tfrac mq(\nu^*(qL)-F)+\{\tfrac mqF\}$.  The divisor $P+\Delta$ is
integral, because $-\lfloor\tfrac mqF\rfloor$, $m\nu^*L$ and
$(d-i)\nu^*H$ are; and $\Delta$ is an effective simple normal crossings
$\mathbf Q$-divisor with $\lfloor\Delta\rfloor=0$.

\emph{Step 3: positivity of $P$.}  By
\Cref{lem:ch40-J-basic}\,\cref{itm:ch40-J:coherent} the sheaf $J_q(qL)$ is
generated relative to $\mathbf A^1$; pulling back to $T$, the line
bundle $\nu^*(qL)-F$ is generated relative to $\mathbf A^1$, hence
$\rho\nu$-semiample and in particular $\rho\nu$-nef.  The bundle
$\nu^*H$ is $\rho\nu$-nef and $\rho\nu$-big, being the pullback of a
$\rho$-ample bundle along a birational morphism, and $d-i\ge d-n\ge1$.
Hence $P$ is $\rho\nu$-nef and $\rho\nu$-big.

\emph{Step 4: vanishing and regularity.}  Relative Kawamata--Viehweg
vanishing, \cite[\S4.3]{PAG1} in its relative form over the affine base
$\mathbf A^1$, applied to the integral divisor $P+\Delta$ with
$P$ nef and big and $\Delta$ a simple normal crossings boundary with
$\lfloor\Delta\rfloor=0$, gives
\[
 R^j(\rho\nu)_*\cO_T(K_T+P+\Delta)=0
 \qquad(j\ge1).
\]
Combining with Step~1, $R^i\rho_*(G(-iH))=0$ for $1\le i\le n$; for
$i>n$ the same vanishing holds because the fibres of $\rho$ have
dimension $n$.  Thus $G$ is $0$-regular relative to $\rho$ with
respect to $H$, and relative Castelnuovo--Mumford regularity,
\cite[\S1.8]{PAG1} applied over the noetherian base ring $k[\varpi]$,
shows that $G=B_m(mL+A)$ is generated relative to $\mathbf A^1$.

\emph{Step 5: the Fubini--Study conclusion.}  The ideal $B_m$ is nonzero
and $\mathbf G_m$-equivariant by
\Cref{prop:ch40-engine-package}\,\cref{itm:ch40-eng:contain} and
\cref{itm:ch40-eng:equivariant}, and $B_m(m\cdot(L+A/m))=B_m(mL+A)$ is
generated relative to $\mathbf A^1$.  Now
\Cref{lem:ch40-equivariant-ideal-fs}\,\cref{itm:ch40-eqfs:fs}, applied with
$N=L+A/m$ and the degree $m$, gives
$u_m=m^{-1}\phi_{B_m}\in\cH^{\mathrm{gf}}_{\mathbf Q}(L+A/m)$.
\end{proof}

\begin{remark}[What Step~3 needs, and why the log resolution is harmless
here]
\label{rem:ch40-char-zero-resolution}
The log resolution $\nu$ is not assumed equivariant, and it is not: the
resolution algorithm need not respect the $\mathbf G_m$-action, and even
if it did, the argument would not use it.  Equivariance of $B_m$ is proved
independently in
\Cref{prop:ch40-engine-package}\,\cref{itm:ch40-eng:equivariant} by a
Vandermonde argument, precisely so that the geometric input and the
representation-theoretic input can be separated.  This separation is what
makes the positive-characteristic proof below structurally parallel: there
is no resolution at all, and the equivariance argument is unchanged.
\end{remark}

\subsection{Positive characteristic: relative generation by Frobenius}
\label{sec:ch40-generation-char-p}

\begin{proposition}\label{prop:bj-smooth-uniform-generation}
For every $m>0$, the equivariant ideal $B_m$ satisfies
\[
                         B_m(mL+A)
\]
generated relative to $\mathbf A^1$.  Consequently,
\[
                         u_m\in\cH^{\mathrm{gf}}_{\mathbf Q}(L+A/m).
\tag{40.13}\label{eq:bj-smooth-gf-class}
\]
The line bundle $A$ is independent of $m$.
\end{proposition}

\begin{proof}
Fix $m$.  For a sufficiently divisible $q$, asymptotic stabilization
gives
\[
                         B_m=\tau(J_q^{m/q}).
\]
For all large Frobenius exponents $e$, set
\[
                         r_e:=\left\lceil\frac{mp^e}{q}\right\rceil.
\]
The stabilized root construction is the surjection
\[
 F_*^e(J_q^{r_e}\otimes\omega_S)
       \twoheadrightarrow B_m\otimes\omega_S.
\tag{40.14}\label{eq:bj-smooth-trace-quotient}
\]
Here $\omega_S$ is the dualizing sheaf, which for the smooth $S$ is the
line bundle $\cO_S(K_S)$; the map is the Grothendieck trace
$\operatorname{Tr}^e$ of \eqref{eq:appC-test-trace}, and
\eqref{eq:bj-smooth-trace-quotient} holds for $e\gg0$ because the sum in
\eqref{eq:appC-test-trace} is increasing in $e$ and stabilizes by
noetherianity, so that a single term realizes it.

Relative generation of $J_q^{r_e}(qr_eL)$ gives a finite free relative
evaluation.  In detail: $J_q(qL)$ is generated relative to
$\mathbf A^1$ by \Cref{lem:ch40-J-basic}\,\cref{itm:ch40-J:coherent}, hence
so is $J_q^{r_e}(qr_eL)$; as $\mathbf A^1$ is affine and
$\rho_*\bigl(J_q^{r_e}(qr_eL)\bigr)$ is a finitely generated
$k[\varpi]$-module, there is a surjection
\[
 \cO_S^{\oplus N_e}\twoheadrightarrow J_q^{r_e}(qr_eL)
\]
for some $N_e\ge1$.
Tensor it by $\cO_S(K_S-qr_eL)$, apply the exact finite
pushforward $F_*^e$, compose with
\eqref{eq:bj-smooth-trace-quotient}, and tensor the target by
\[
                         mL+A-K_S=mL+dH.
\]
The Frobenius projection formula gives a surjection from a finite direct sum
of copies of
\[
 F_*^e\cO_S\bigl(
 K_S+(mp^e-qr_e)L+p^edH
 \bigr)
 \longrightarrow B_m(mL+A).
\tag{40.15}\label{eq:bj-smooth-root-source}
\]
The twist is exact because
\[
 K_S=\operatorname{pr}_X^*K_X,\qquad
 K_{\mathbf A^1}\simeq\cO_{\mathbf A^1}.
\tag{40.16}\label{eq:bj-smooth-canonical-product}
\]
The computation behind \eqref{eq:bj-smooth-root-source} is the identity
\[
 \begin{aligned}
 F^e_*&\bigl(\cO_S(K_S-qr_eL)\bigr)\otimes\cO_S(mL+dH)\\
 &=F^e_*\bigl(\cO_S(K_S-qr_eL)\otimes(F^e)^*\cO_S(mL+dH)\bigr)\\
 &=F^e_*\bigl(\cO_S(K_S+(mp^e-qr_e)L+p^edH)\bigr),
 \end{aligned}
\]
using $(F^e)^*M=M^{\otimes p^e}$ for a line bundle $M$.

For $1\le i\le n$, twist the source in
\eqref{eq:bj-smooth-root-source} by $-iH$.  Its $i$-th higher direct
image under $\rho$ is
\[
 F_{\mathbf A^1,*}^eR^i\rho_*
 \cO_S\bigl(
 K_S+(mp^e-qr_e)L+p^e(d-i)H
 \bigr).
\tag{40.17}\label{eq:bj-smooth-root-cohomology}
\]
Two identities are used here.  The projection formula again converts the
twist by $-iH$ into the twist by $-p^eiH$ inside the Frobenius
pushforward; and $\rho\circ F_S^e=F_{\mathbf A^1}^e\circ\rho$, with
$F^e_{\mathbf A^1}$ affine and hence exact, converts
$R^i\rho_*F^e_{S,*}$ into $F^e_{\mathbf A^1,*}R^i\rho_*$.

The residual coefficient belongs to the fixed finite set
\[
 mp^e-q\left\lceil\frac{mp^e}{q}\right\rceil
     \in\{-(q-1),\ldots,-1,0\}.
\tag{40.18}\label{eq:bj-smooth-residual-coefficient}
\]
For a coherent sheaf $E$ on $X$,
\[
 R^i\rho_*\operatorname{pr}_X^*E
       =H^i(X,E)\otimes_k k[\varpi].
\tag{40.19}\label{eq:bj-smooth-product-cohomology}
\]
Since $p^e(d-i)\to+\infty$ for $i\le n$, Serre vanishing makes
\eqref{eq:bj-smooth-root-cohomology} zero for $e\gg0$; for $i>n$,
it vanishes by dimension.  Relative Castelnuovo--Mumford regularity makes
\eqref{eq:bj-smooth-root-source} relatively generated, and so is its
quotient $B_m(mL+A)$.

Two points in the last paragraph should be made explicit, because they are
where the uniformity in $m$ resides.  First, by
\eqref{eq:bj-smooth-residual-coefficient} the bundle
$K_X+cL$ occurring in \eqref{eq:bj-smooth-root-cohomology} ranges over the
\emph{finite} set indexed by $c\in\{-(q-1),\ldots,0\}$, so a single Serre
threshold $N_0=N_0(m,q)$ serves for all of them: there is $N_0$ with
$H^i(X,K_X+cL+NH)=0$ for all $i\ge1$, all such $c$, and all
$N\ge N_0$.  Since $d-i\ge1$, the coefficient $p^e(d-i)\ge p^e$
exceeds $N_0$ for $e$ large, and $e$ may be chosen as large as we
please, subject only to \eqref{eq:bj-smooth-trace-quotient}, which also
holds for all large $e$.  Second, the threshold $N_0$ depends on $m$
and $q$, but $A$ does not: the vanishing is used to produce
\emph{some} surjection onto $B_m(mL+A)$ for each $m$ separately, and the
only object required to be independent of $m$ is the twist $A$, which
was fixed in \eqref{eq:bj-smooth-fixed-twist} before $m$ was chosen.

Equivariance of $B_m$ follows because scaling automorphisms preserve
$J_\bullet$ and test ideals commute with isomorphisms; this is
\Cref{prop:ch40-engine-package}\,\cref{itm:ch40-eng:equivariant}.
Homogeneous relative generators give \eqref{eq:bj-smooth-gf-class}: apply
\Cref{lem:ch40-equivariant-ideal-fs}\,\cref{itm:ch40-eqfs:fs} with
$N=L+A/m$ and degree $m$, exactly as in Step~5 of the proof of
\Cref{prop:ch40-multiplier-uniform-generation}.
\end{proof}

\subsection{The drift, and why it is oriented}
\label{sec:ch40-drift}

From now on we enlarge $d$ once, if necessary, so that in addition to
$d\ge n+1$ the class $c_1(A)=c_1(K_X)+d\,c_1(H)$ is \emph{nef}; this
is possible because $H$ is ample.  Nefness, not ampleness, is what
\Cref{lem:ch40-drift-nef-monotone} needs, and the weaker demand is the
right one: on $X=\mathbf P^n$ with $H=\cO(1)$ and $d=n+1$ one has
$A=\cO_X$ and $c_1(A)=0$, which is nef but not ample, and that is
exactly the case computed in \Cref{ex:ch40-p2-engines} and
\Cref{ex:ch40-drift-real}\,(a).  The enlargement does not affect
\Cref{prop:ch40-multiplier-uniform-generation} or
\Cref{prop:bj-smooth-uniform-generation}, whose hypotheses only bound $d$
from below.  Write
\[
 \theta_m:=c_1(L)+\frac1m c_1(A)\in N^1(X)_{\mathbf Q},
 \qquad m\ge1 .
\]

\begin{lemma}
\label{lem:ch40-drift-nef-monotone}
For $m\mid m'$ one has $\theta_m-\theta_{m'}
=\bigl(\tfrac1m-\tfrac1{m'}\bigr)c_1(A)\in\Nef(X)$, and
$\theta_m\to c_1(L)$ along divisibility.  Consequently:
\begin{enumerate}
\item\label{itm:ch40-drift:tail}
 for every $m_0$ and every multiple $m$ of $m_0$,
 $u_m\in\PSH_X(\theta_{m_0})$;
\item\label{itm:ch40-drift:decreasing}
 the tail $\{u_m:m_0\mid m\}$ is a decreasing net in the single class
 $\PSH_X(\theta_{m_0})$.
\end{enumerate}
\end{lemma}

\begin{proof}
The class $\bigl(\tfrac1m-\tfrac1{m'}\bigr)c_1(A)$ is a nonnegative
multiple of a nef class, hence nef, and
$\lVert\theta_m-c_1(L)\rVert=m^{-1}\lVert c_1(A)\rVert\to0$.

\cref{itm:ch40-drift:tail}  By
\Cref{prop:ch40-multiplier-uniform-generation}, respectively
\Cref{prop:bj-smooth-uniform-generation},
$u_m\in\cH^{\mathrm{gf}}_{\mathbf Q}(L+A/m)\subseteq\PSH_X(\theta_m)$, the
inclusion being \Cref{def:bj-psh} applied to the constant net.  Since
$\theta_{m_0}-\theta_m$ is nef, \Cref{lem:ch34-psh-nonempty} gives
$0\in\PSH_X(\theta_{m_0}-\theta_m)$ and
\Cref{thm:bj-basic-closure}\,(3) gives $u_m\in\PSH_X(\theta_{m_0})$.

\cref{itm:ch40-drift:decreasing}  Decreasingness is
\Cref{prop:ch40-engine-package}\,\cref{itm:ch40-eng:sandwich}.
\end{proof}

\begin{warning}{The class error is exact, not an estimate}
\label{warn:ch40-exact-class}
The class carried by $u_m$ is \emph{exactly} $\theta_m=c_1(L)+m^{-1}c_1(A)$,
not merely a class within an unspecified $O(m^{-1})$ neighbourhood of
$c_1(L)$.  This matters twice.

It matters for the limit.  \Cref{warn:ch38-drift-not-absorbed} records that
``the error tends to zero, so it may be ignored'' is not an argument, and
that the correct hypotheses are order hypotheses in the nef cone.  Here the
error is oriented, by \Cref{lem:ch40-drift-nef-monotone}, and in the
same nef-cone direction as in \Cref{sec:ch40-real-to-rational}: the stage
classes \emph{dominate} the target and decrease to it.  Unlike the rational
approximation there, they also form a monotone chain.  This is exactly the
configuration that clause~(5) of \Cref{thm:bj-basic-closure} was stated for,
and \Cref{lem:ch40-drift-nef-monotone}\,\cref{itm:ch40-drift:tail} converts it
into the configuration of clause~(4) inside each fixed class
$\theta_{m_0}$, so that either clause may be used.  Both routes are
written out in \Cref{sec:ch40-psh-limit}.

It matters for bookkeeping.  If one only knew
$\lVert\theta_m-c_1(L)\rVert=O(m^{-1})$, one could not conclude that the
tail lies in a fixed class, and the two-step argument of
\Cref{lem:ch40-drift-nef-monotone} would be unavailable; one would be forced
to use clause~(5) with no fallback.  Knowing the class exactly costs
nothing here and is worth the notation.
\end{warning}

\begin{example}[The two $O(m^{-1})$ errors are real, and independent]
\label{ex:ch40-drift-real}
Two errors of size $O(m^{-1})$ occur in this chapter: the class error
$\theta_m-c_1(L)=m^{-1}c_1(A)$, an element of $N^1(X)_{\mathbf Q}$, and
the value error $u_m-Q$, a nonnegative function on $X^{\An}$.  Neither
is an artefact, and neither controls the other.

\emph{(a) The class error can vanish while the value error does not.}  Let
$X=\mathbf P^n$, $H=\cO(1)$, $d=n+1$.  Then
\[
 A=K_X+(n+1)H=\cO(-(n+1))+\cO(n+1)=\cO_X,
\]
so $\theta_m=c_1(L)$ for every $m$ and the class error is identically
zero.  Nevertheless, in \Cref{ex:ch40-p2-engines} with $n=2$,
\[
 u_m-Q=\frac2m\min\{s,1\}
 \qquad(m\ge2),
\]
which is strictly positive at every point with $s>0$ and for every finite
$m$.  Thus no finite stage computes the envelope, even when the classes are
exactly right.

\emph{(b) The class error is nonzero for every admissible choice of the
auxiliary data, on most varieties.}  Let $X$ be a smooth projective curve
of genus $g\ge1$, so $n=1$.  For any very ample globally generated
$H$ and any $d\ge2$,
\[
 \deg A=(2g-2)+d\deg H\ \ge\ 0+2\cdot1=2>0,
\]
so $c_1(A)\ne0$ and $\theta_m\ne c_1(L)$ for \emph{every} $m$ and
\emph{every} admissible choice of $H$ and $d$.  Concretely, for $E$
elliptic, $P\in E$, $H:=\cO_E(3P)$ and $d=2$ one gets
$A=K_E+2H=\cO_E(6P)$ of degree $6$, so that
$\deg\theta_m=\deg L+6/m$.  On such an $X$ the construction never
certifies the class $c_1(L)$ at a finite stage: what
\Cref{prop:ch40-multiplier-uniform-generation} and
\Cref{prop:bj-smooth-uniform-generation} produce are sections of
$mL+A$, and the class $c_1(L)$ is reached only in the limit.  The same
computation on $X=\mathbf P^1$, where $K_X+2H=\cO$ for
$H=\cO(1)$ and $d=2$, gives $\deg A=0$; genus zero is the exceptional
case, exactly as in (a).

\emph{(c) Conclusion.}  The two errors live in different spaces, vanish
under different circumstances, and are removed by different devices: the
class error by the limit clauses of \Cref{thm:bj-basic-closure}, the value
error by the valuation inequality of \Cref{sec:ch40-valuation}.  This is the
distinction drawn in \Cref{rem:ch38-two-errors} in the abstract; here it is
instantiated with numbers.
\end{example}

\begin{remark}[Two features of the twist that are not accidental]
\label{rem:ch40-twist-features}
The reader may ask why the argument produces $K_X+dH$ rather than a twist
adapted to $L$ and the obstacle.  Two reasons, one in each characteristic.

In characteristic zero, the sheaf whose higher direct images are computed is
$\cO_T(K_T+P+\Delta)$, and the canonical bundle is forced: the vanishing
theorem is a statement about $K_T+\text{nef and big}$.  The passage from
$K_T$ to $K_X$ is the relative canonical formula
\eqref{eq:bj-smooth-canonical-product} together with
$K_{T/S}$ being absorbed into the multiplier ideal.

In characteristic $p$, the sheaf that surjects onto $B_m\otimes\omega_S$
is a Frobenius pushforward of a twist of $\omega_S$, by the very
definition \eqref{eq:appC-test-trace} of the test ideal, and again the
canonical bundle is forced.  What is \emph{not} forced is that the residual
coefficient \eqref{eq:bj-smooth-residual-coefficient} stays in a bounded set
while the ample coefficient grows like $p^e$; this is the specifically
positive-characteristic phenomenon, and it is the reason a twist independent
of $m$ exists at all.  If $r_e$ were replaced by $mp^e/q$ without
rounding, there would be no residual term and no need for the boundedness
observation --- but $mp^e/q$ is not an integer, and the rounding is
exactly what the trace description requires.
\end{remark}

\section{The valuation inequality and the psh limit}
\label{sec:ch40-valuation}

The single valuation estimate needed below is worth isolating.  It is the
place where positive characteristic enters the identification of the
limit.

\subsection{The estimate in positive characteristic}
\label{sec:ch40-canton}

\begin{lemma}\label{lem:bj-smooth-canton}
Let $I$ be a nonzero coherent ideal on a regular $F$-finite scheme
$S$, let $t\ge0$, and let $w$ be a divisorial valuation centered on
$S$.  Then
\[
                       w\bigl(\tau(I^t)\bigr)
                          \ge t\,w(I)-A_S(w).
\tag{40.20}\label{eq:bj-smooth-canton-test}
\]
\end{lemma}

\begin{proof}
Choose a local generator of $\omega_S$ near the center of $w$.  Under
this trivialization the Grothendieck trace becomes a $p^{-e}$-linear map
\[
                         \Phi_e:F_*^e\cO_S\longrightarrow\cO_S.
\]
Such a trivialization exists because $S$ is regular, so $\omega_S$ is
invertible; and $\Phi_e$ generates
$\operatorname{Hom}_{\cO_S}(F^e_*\cO_S,\cO_S)$ as an
$F^e_*\cO_S$-module, by Kunz's criterion and Grothendieck duality for the
finite flat morphism $F^e$.  Consequently the trace description
\eqref{eq:appC-test-trace} reads, after the trivialization,
\[
 \tau(I^t)=\sum_{e\ge0}\Phi_e\bigl(F^e_*I^{\lceil tp^e\rceil}\bigr),
\]
a sum which is increasing in $e$ and stabilizes.

If $g=\Phi_e(f)\ne0$, Canton's definition and the identification of the
Cartier-algebra discrepancy with the ordinary discrepancy of the regular
pair $(S,0)$ give
\[
 \frac{w(f)-p^e w(g)}{p^e-1}
     \le A_S(w).
\tag{40.21}\label{eq:bj-smooth-canton-trace}
\]
The root construction of $\tau(I^t)$ uses elements
$f\in I^{\lceil tp^e\rceil}$.  Therefore every trace generator $g$
satisfies
\[
 w(g)\ge
 \frac{\lceil tp^e\rceil}{p^e}w(I)
 -\left(1-\frac1{p^e}\right)A_S(w).
\]
Indeed, solving \eqref{eq:bj-smooth-canton-trace} for $w(g)$ gives
$w(g)\ge p^{-e}w(f)-(1-p^{-e})A_S(w)$, and
$w(f)\ge\lceil tp^e\rceil\,w(I)$ because
$f\in I^{\lceil tp^e\rceil}$.  Since
$\lceil tp^e\rceil\ge tp^e$ and $w(I)\ge0$, the right-hand side is at
least $t\,w(I)-A_S(w)$.

Take the minimum over the trace generators and then the stabilized limit.
An arbitrary local section of $\tau(I^t)$ is an $\cO_S$-linear
combination of such $g$, and $w$ of a sum is at least the minimum of the
values, so the bound passes from the generators to the ideal.  This proves
\eqref{eq:bj-smooth-canton-test}; see
\cite[Definition~3.2 and Proposition~4.1]{Canton}.
\end{proof}

\begin{remark}[The exact shape of the characteristic-$p$ input]
\label{rem:ch40-canton-shape}
Three features of \eqref{eq:bj-smooth-canton-trace} are worth naming.

It is an inequality about \emph{one} Frobenius level $e$, with the factor
$(p^e-1)^{-1}$, and it is uniform in $e$; this is what allows the
stabilized limit to be taken without losing the bound.  Its right-hand side
is the ordinary log discrepancy $A_S(w)$ of \eqref{eq:appC-log-discrepancy},
not a characteristic-dependent substitute: Canton's theorem is precisely the
identification of the Cartier-algebra invariant, defined by
\eqref{eq:bj-smooth-canton-trace}, with the birational invariant.  And it is
a statement about a \emph{fixed} $w$; no bound uniform over all divisorial
valuations is asserted, and none is available, since $A_S$ is unbounded on
$S^{\divis}$.

The last point is what dictates the shape of
\Cref{prop:bj-smooth-divisorial-identification}: the limit is identified
pointwise on $X^{\divis}$, one valuation at a time, and the passage from
divisorial equality to equality on $X^{\An}$ is then made by the psh
comparison \Cref{prop:ch39-determination-package}, not by a uniform estimate.
\end{remark}

\subsection{The estimate in characteristic zero}
\label{sec:ch40-multiplier-valuation}

The characteristic-zero counterpart of \Cref{lem:bj-smooth-canton} is
\eqref{eq:appC-multiplier-valuation-bound}, and its proof is one line from
the valuative description \eqref{eq:appC-valuative-multiplier}: if
$f\in\mathcal J_S(I^t)$ then $v(f)+A_S(v)>t\,v(I)$ for every divisorial
$v$, so $v(\mathcal J_S(I^t))\ge t\,v(I)-A_S(v)$, the strict inequality
being weakened to a non-strict one.  We record the parallel and the
difference.

\begin{lemma}
\label{lem:ch40-valuation-both}
Let $w$ be a divisorial valuation centred on $S$.  Then, with
$B_m$ as in \Cref{def:ch40-stage-ideals} and
$w(J_\bullet)=\inf_qq^{-1}w(J_q)$,
\[
 w(B_m)\ \ge\ m\,w(J_\bullet)-A_S(w)
\]
in both characteristics; and in both characteristics
$w(B_m)\le w(J_m)$, because $J_m\subseteq B_m$.  The same constant
$A_S(w)$, the ordinary log discrepancy of the regular scheme $S$,
appears in both.
\end{lemma}

\begin{proof}
This is \Cref{prop:ch40-engine-package}\,\cref{itm:ch40-eng:valuation} and
\cref{itm:ch40-eng:contain}; the characteristic-zero half is
\eqref{eq:appC-asymptotic-discrepancy} and the positive-characteristic half
is \Cref{lem:bj-smooth-canton} applied at a sufficiently divisible $q$.
\end{proof}

The identity of the two constants is a coincidence of statements, not of
proofs.  In characteristic zero $A_S(w)$ enters as the coefficient of the
relative canonical divisor on a log resolution, through the floor function
in \eqref{eq:appC-multiplier-resolution}; in characteristic $p$ it enters
as the Cartier-algebra discrepancy of \eqref{eq:bj-smooth-canton-trace}, and
the theorem that the two agree is Canton's.  Neither proof is available in
the other characteristic.

\subsection{The Gauss extension of a divisorial point}
\label{sec:ch40-gauss-divisorial}

To apply \Cref{lem:ch40-valuation-both} at a point $v\in X^{\divis}$ we
need to know that its Gauss extension is divisorial on $S$ and to know its
log discrepancy.  Both are explicit.

\begin{lemma}
\label{lem:ch40-gauss-log-discrepancy}
Let $v\in X^{\divis}$.  Then $w:=\sigma(v)$ is a divisorial valuation on
$S$, centred on the central fibre $X\times\{0\}$, and
\[
 A_S\bigl(\sigma(v)\bigr)=A_X(v)+1,
\]
where $A_X$ is the log discrepancy \eqref{eq:appC-log-discrepancy},
extended homogeneously to positive rational multiples of divisorial
valuations.  In particular $A_S(\sigma(v))<\infty$ for every
$v\in X^{\divis}$, and $A_S(\sigma(v_{\triv}))=1$.
\end{lemma}

\begin{proof}
If $v=v_{\triv}$, then $\sigma(v)(\sum_if_i\varpi^i)=\min\{i:f_i\ne0\}
=\ord_{X\times\{0\}}$, a divisorial valuation with
$A_S=1$, and $A_X(v_{\triv})=0$; so the formula holds.

Otherwise $v=t\ord_F$ with $t\in\mathbf Q_{>0}$ and $F$ a prime
divisor on a smooth projective model $\pi:Y\to X$.  On
$Y\times\mathbf A^1$ the two divisors $F\times\mathbf A^1$ and
$Y\times\{0\}$ are smooth and meet transversally, so
$(Y\times\mathbf A^1,\ F\times\mathbf A^1+Y\times\{0\})$ is a simple
normal crossings pair in the sense of \Cref{def:ch32-snc-pair}.  Directly
from \eqref{eq:ch40-gauss}, the semivaluation $\sigma(t\ord_F)$ is the
quasi-monomial valuation of that pair with weights $(t,1)$.  Since the
weights are rational, $\sigma(v)$ is divisorial by
\Cref{thm:ch32-quasimonomial-invariants}\,\cref{itm:ch32-qi:divisorial}, and
by \Cref{prop:ch32-log-discrepancy-affine} its log discrepancy relative to
the snc pair is $t+1$.

It remains to convert this into a log discrepancy over $S$.  Let
$\Pi:=\pi\times\operatorname{id}:Y\times\mathbf A^1\to S$.  Then
$K_{(Y\times\mathbf A^1)/S}=\operatorname{pr}_Y^*K_{Y/X}$, so
\[
 \begin{aligned}
 A_S(\sigma(v))
 &=A_{Y\times\mathbf A^1}(\sigma(v))
   +\sigma(v)\bigl(K_{(Y\times\mathbf A^1)/S}\bigr)\\
 &=(t+1)+t\,\ord_F(K_{Y/X})
 =t\bigl(1+\ord_F(K_{Y/X})\bigr)+1,
 \end{aligned}
\]
which is $A_X(v)+1$ by \eqref{eq:appC-log-discrepancy} and homogeneity.
\end{proof}

\begin{example}[The valuation inequality checked by hand on a monomial
model]
\label{ex:ch40-monomial-valuation}
Return to \Cref{ex:ch40-p2-obstacle} and
\Cref{ex:ch40-p2-engines}, and work in the regular local ring at
$(p,0)\in S$ with regular system of parameters $(x,y,\varpi)$.  For
$a,b,c\in\mathbf Q_{>0}$ let $w_{a,b,c}$ be the monomial valuation with
$w(x)=a$, $w(y)=b$, $w(\varpi)=c$; it is quasi-monomial for the snc
divisor $\{xy\varpi=0\}$, divisorial when the weights are rational, and
\[
 A_S(w_{a,b,c})=a+b+c
\]
by \Cref{prop:ch32-log-discrepancy-affine}.  We have
\[
 \begin{gathered}
 w(J_\bullet)=w(J_1)=w(\varpi\mathfrak m)=c+\min\{a,b,c\},\\
 w(B_m)=w\bigl(\varpi^m\mathfrak m^{m-2}\bigr)=mc+(m-2)\min\{a,b,c\}
 \quad(m\ge2).
 \end{gathered}
\]
The inequality of \Cref{lem:ch40-valuation-both} therefore reads
\[
 mc+(m-2)\min\{a,b,c\}
 \ \ge\ m\bigl(c+\min\{a,b,c\}\bigr)-(a+b+c),
\]
that is,
\[
 2\min\{a,b,c\}\ \le\ a+b+c,
\]
which is true, and strictly so for positive weights: the deficit is
$a+b+c-2\min\{a,b,c\}>0$.  Thus the estimate holds, is independent of
$m$, and is \emph{not} sharp; the true error $2\min\{a,b,c\}$ is smaller
than the guaranteed error $A_S(w)=a+b+c$.  Both are constants, and both
disappear after division by $m$, which is all the argument needs.

Specializing to a Gauss extension: for $v=t\ord_p\in X^{\divis}$ with
$t\in\mathbf Q_{>0}$, one has $\sigma(v)=w_{t,t,1}$, and
\Cref{lem:ch40-gauss-log-discrepancy} gives
$A_S(\sigma(v))=A_X(t\ord_p)+1=2t+1$, in agreement with
$a+b+c=t+t+1$.  The inequality becomes
$2\min\{t,1\}\le2t+1$, true for every $t>0$.  Dividing by $m$ and
negating, the resulting bound on functions is
\[
 Q(v)\ \le\ u_m(v)\ \le\ Q(v)+\frac{2t+1}{m},
\]
whereas the exact value computed in \Cref{ex:ch40-p2-engines} is
$u_m(v)=Q(v)+\tfrac2m\min\{t,1\}$.  The estimate is correct, crude, and
sufficient.
\end{example}

\subsection{The limit is plurisubharmonic}
\label{sec:ch40-psh-limit}

Restrict henceforth to a cofinal factorial sequence of indices, so that
the $u_m$ decrease.  Fix one index $m_0$.  The GF function
\[
                         q_{m_0}:=\frac1{m_0}\phi_{J_{m_0}}
\]
is generically finite.  The complete base-ideal construction gives
\[
                         q_{m_0}\le Q\le u_m
\tag{40.22}\label{eq:bj-smooth-gf-lower-bound}
\]
for every sufficiently divisible $m$.  Thus
\[
                             u:=\lim_m u_m
\tag{40.23}\label{eq:bj-smooth-limit}
\]
is generically finite.  Since $c_1(L+A/m)\to c_1(L)$,
\eqref{eq:bj-smooth-gf-class} and the decreasing-limit axiom give
\[
                              u\in\PSH_X(L).
\tag{40.24}\label{eq:bj-smooth-limit-psh}
\]
The explicit lower bound in \eqref{eq:bj-smooth-gf-lower-bound} rules out
the identically $-\infty$ alternative.  This is important: the class
convergence alone does not do so.

We spell out the three assertions.

\emph{Generic finiteness.}  The function $q_{m_0}$ is a generically finite
Fubini--Study function by
\Cref{prop:ch40-Q-limit}\,\cref{itm:ch40-Q:gf}, hence finite on the
full-support locus $X^{\val}$ by
\Cref{prop:ch34-generic-finiteness}\,\cref{itm:ch34-gf:equiv}.  The first
inequality of \eqref{eq:bj-smooth-gf-lower-bound} is
\Cref{prop:ch40-Q-limit}\,\cref{itm:ch40-Q:sup}, and the second is
\Cref{prop:ch40-engine-package}\,\cref{itm:ch40-eng:sandwich}.  Therefore
$u\ge q_{m_0}>-\infty$ on $X^{\val}$, and in particular
$u(v_{\triv})>-\infty$; by
\Cref{prop:ch34-generic-finiteness}\,\cref{itm:ch34-gf:equiv} again, $u$ is
generically finite.  This is the step that fails if one tries to run the
construction with $J_\bullet=0$, and it is the reason the Veronese of
\Cref{lem:ch40-J-basic}\,\cref{itm:ch40-J:nonzero} was performed.

\emph{First route to \eqref{eq:bj-smooth-limit-psh}: varying classes.}  The
net $(u_m)$, indexed by the divisibility order, is decreasing by
\Cref{prop:ch40-engine-package}\,\cref{itm:ch40-eng:sandwich}, consists of
generically finite Fubini--Study functions for the $\mathbf Q$-line
bundles $L+A/m$ by \eqref{eq:bj-smooth-gf-class}, and satisfies
$c_1(L+A/m)\to c_1(L)$.  Its pointwise limit $u$ is generically finite.
These are exactly the three conditions of \Cref{def:bj-psh}, so
$u\in\PSH_X(c_1(L))$.  Upper semicontinuity is automatic, an infimum of
continuous functions being upper semicontinuous.

\emph{Second route: a fixed class, then a limit of classes.}  Fix $m_0$.
By \Cref{lem:ch40-drift-nef-monotone}\,\cref{itm:ch40-drift:tail}, every
$u_m$ with $m_0\mid m$ lies in the single space
$\PSH_X(\theta_{m_0})$, and the tail is decreasing with generically finite
limit $u$; so \Cref{thm:bj-basic-closure}\,(4) gives
$u\in\PSH_X(\theta_{m_0})$.  As $m_0$ was arbitrary and
$\theta_{m_0}\to c_1(L)$, the closedness assertion of
\Cref{thm:bj-basic-closure}\,(5) gives $u\in\PSH_X(c_1(L))$.

Both routes are recorded because they consume different clauses.  The first
uses \Cref{def:bj-psh} directly and is shorter; the second never leaves a
fixed class until the very last step and is the one that matches the drift
discipline of \Cref{warn:ch38-drift-not-absorbed}.  Neither uses the
envelope property, closure under suprema, or
\Cref{thm:ch34-fs-regularization}.

\begin{warning}{Generic finiteness is not automatic}
\label{warn:ch40-generic-finiteness}
A decreasing net of generically finite Fubini--Study functions in a
\emph{fixed} class can converge to $-\infty$ identically.  For the
simplest instance, let $L$ be a line bundle with a nonzero section $s$
and put $\phi_m:=\log|s|-m$; each $\phi_m$ lies in
$\cH^{\mathrm{gf}}_{\mathbf R}(L)$, being the one-element weighted section
system $(1,(s),(-m))$, the net is decreasing, the classes are constant, and
$\lim_m\phi_m\equiv-\infty$.  \Cref{def:bj-psh}(2) excludes such limits by
fiat, so a proof that produces a decreasing net must independently produce a
lower bound.

Here the lower bound is \eqref{eq:bj-smooth-gf-lower-bound}, and two of its
features are used.  It is a bound by a \emph{function} and not by a
constant, which is what is needed, since $u$ may legitimately be
$-\infty$ on the common base locus and no constant bound could hold there;
and the bounding function $q_{m_0}$ is itself generically finite, so that
the bound is nonvacuous at $v_{\triv}$.  \Cref{ex:ch34-decreasing-to-pole}
exhibits the subtler phenomenon, in which the limit is finite somewhere and
$-\infty$ on a nonempty proper analytic subset; that behaviour is
permitted, and indeed occurs here whenever $L$ has base points.
\end{warning}

\subsection{Identification of the divisorial values}
\label{sec:ch40-divisorial-identification}

\begin{proposition}\label{prop:bj-smooth-divisorial-identification}
For every divisorial point $v\in X^{\divis}$, including the trivial
valuation,
\[
                              u(v)=Q(v).
\tag{40.25}\label{eq:bj-smooth-divisorial-identification}
\]
\end{proposition}

\begin{proof}
Let $w=\sigma(v)$ be the Gauss extension to $S$, normalized by
$w(\varpi)=1$.  If $v$ is a rational multiple of a divisorial valuation,
then $w$ is the corresponding rational toroidal valuation.  If $v$ is
trivial, then $w=\ord_{X\times\{0\}}$.  In either case $w$ is
divisorial and $A_S(w)<\infty$; this is
\Cref{lem:ch40-gauss-log-discrepancy}, which moreover computes
$A_S(w)=A_X(v)+1$.

Set
\[
                         \alpha_w:=\inf_q\frac{w(J_q)}q.
\tag{40.26}\label{eq:bj-smooth-alpha}
\]
For a sufficiently divisible $q$, asymptotic stabilization and
\Cref{lem:bj-smooth-canton}, applied to $I=J_q$ and $t=m/q$, give
\[
                         w(B_m)\ge\frac mqw(J_q)-A_S(w).
\]
In characteristic zero the same inequality is
\eqref{eq:appC-asymptotic-discrepancy}; the two cases are the two halves of
\Cref{lem:ch40-valuation-both}, and the argument does not distinguish them
from here on.
The inclusion $J_m\subset B_m$ supplies the other inequality.  Hence
\[
              m\alpha_w-A_S(w)\le w(B_m)\le w(J_m).
\tag{40.27}\label{eq:bj-smooth-canton-squeeze}
\]
The limit defining $\alpha_w$ exists along divisibility because
\eqref{eq:bj-smooth-graded-ideals} implies
\[
                     w(J_{m+r})\le w(J_m)+w(J_r).
\]
Divide \eqref{eq:bj-smooth-canton-squeeze} by $m$, negate, and use the
base-ideal limit \eqref{eq:bj-smooth-Q}.  The discrepancy error is exactly
$A_S(w)/m$, so it disappears for each fixed $w$, proving
\eqref{eq:bj-smooth-divisorial-identification}.  No discrepancy bound
uniform in $v$ is asserted or needed.

In detail: all three quantities in
\eqref{eq:bj-smooth-canton-squeeze} are finite real numbers, because $w$ is
a valuation of $K(S)$ and $J_m$, $B_m$ are nonzero ideals, so no
extended-real arithmetic occurs.  Dividing by $m$ and negating reverses the
inequalities:
\[
 -\frac{w(J_m)}m\ \le\ -\frac{w(B_m)}m\ \le\ -\alpha_w+\frac{A_S(w)}m,
\]
that is, $q_m(v)\le u_m(v)\le-\alpha_w+A_S(w)/m$.  By Fekete's subadditive
lemma applied to the subadditive sequence $m\mapsto w(J_m)$, the left-hand
side converges to $-\alpha_w$ along divisibility; and by
\Cref{prop:ch40-Q-limit}\,\cref{itm:ch40-Q:sup} its limit is $Q(v)$, so
$Q(v)=-\alpha_w$.  Letting $m\to\infty$ along divisibility in the
displayed inequalities gives
\[
 Q(v)\ \le\ u(v)\ \le\ Q(v)+\lim_m\frac{A_S(w)}m=Q(v),
\]
as required.
\end{proof}

\begin{remark}[Why a pointwise identification suffices]
\label{rem:ch40-no-uniform-discrepancy}
The error term in \eqref{eq:bj-smooth-canton-squeeze} is $A_S(w)/m$, and
$A_S(w)$ is not bounded as $w$ ranges over the Gauss extensions of
divisorial points: by \Cref{lem:ch40-gauss-log-discrepancy} it equals
$A_X(v)+1$, which tends to infinity along $v=t\ord_F$ as $t\to\infty$
and also along sequences of divisors $F$ of increasing log discrepancy.
Consequently the convergence $u_m\to Q$ on $X^{\divis}$ is
\emph{pointwise and not uniform}, and no argument below may use uniformity.
This is why the passage from \eqref{eq:bj-smooth-divisorial-identification}
to a statement on all of $X^{\An}$ is made by the psh comparison
\Cref{prop:ch39-determination-package}\,\cref{itm:ch39-det:compare}, whose
hypotheses are a psh function below an upper-semicontinuous one, rather than
by a limit interchange.  \Cref{ex:ch35-evaluation-discontinuous} shows that
evaluation is not continuous at proper-support points, so a uniformity
argument would in any case not extend the identification off
$X^{\divis}$.
\end{remark}

\section{Identification with the obstacle envelope}
\label{sec:ch40-identification}

\begin{theorem}\label{thm:bj-smooth-envelope}
Let $X$ be a smooth projective variety over the algebraically closed
trivially valued field $k$.  Then every class
$\theta\in N^1(X)_{\mathbf R}$ has the envelope property.
\end{theorem}

\begin{proof}
By \Cref{prop:bj-smooth-obstacle-interface} it suffices to prove that
$P_L(\phi_D)\in\PSH_X(c_1(L))$ for every big line bundle $L$ and every
rational PL obstacle, presented as in \eqref{eq:bj-smooth-graph} with
$D\le0$.  Fix such data and construct, as in
\Crefrange{sec:ch40-base-ideals}{sec:ch40-valuation}, the graded sequence
$J_\bullet$, the stage ideals $B_m$ of \Cref{def:ch40-stage-ideals}, the
stage functions $u_m=m^{-1}\phi_{B_m}$, and their decreasing limit $u$ of
\eqref{eq:bj-smooth-limit}.  The construction is the same in every
characteristic; only the definition of $B_m$ changes, and every property of
$B_m$ used below is one of the four listed in
\Cref{prop:ch40-engine-package}, each of which has been established in both
characteristics.

\emph{Step 1: $u$ is psh.}  This is \eqref{eq:bj-smooth-limit-psh}, proved
in \Cref{sec:ch40-psh-limit} by either of the two routes given there.

\emph{Step 2: $u$ lies below the obstacle.}
By \eqref{eq:bj-smooth-obstacle-identity}, one has $Q\le\phi_D$.
Together with \eqref{eq:bj-smooth-divisorial-identification}, divisorial
comparison against the usc obstacle yields
\[
                              u\le\phi_D
                              \quad\text{on }X^{\An}.
\tag{40.28}\label{eq:bj-smooth-limit-below-obstacle}
\]
Explicitly: $q_m\le\phi_D$ for every $m$ by
\Cref{prop:ch40-Q-limit}\,\cref{itm:ch40-Q:gf}, hence $Q\le\phi_D$ on all
of $X^{\An}$; by \eqref{eq:bj-smooth-divisorial-identification},
$u=Q\le\phi_D$ on $X^{\divis}$; and since $u$ is
$c_1(L)$-psh by Step~1 while $\phi_D$ is continuous, hence upper
semicontinuous,
\Cref{prop:ch39-determination-package}\,\cref{itm:ch39-det:compare} promotes
the inequality from $X^{\divis}$ to $X^{\An}$.
Thus $u$ is a competitor for $P_L(\phi_D)$, and $u\le P_L(\phi_D)$.

\emph{Step 3: every competitor lies below $u$.}
Conversely, let $\tau\in\PSH_X(L)$ satisfy $\tau\le\phi_D$.
By \Cref{prop:ch40-envelope-comparison}\,\cref{itm:ch40-cmp:divisorial},
applied with the continuous obstacle $h=\phi_D$,
\[
 \tau\le Q_L(\phi_D)=Q=u
 \qquad\text{on }X^{\divis},
\]
the last equality being
\eqref{eq:bj-smooth-divisorial-identification}.  Both $\tau$ and $u$ are
psh, so a second application of
\Cref{prop:ch39-determination-package}\,\cref{itm:ch39-det:compare}, with the
psh function $\tau$ below the upper-semicontinuous function $u$, gives
$\tau\le u$ on all of $X^{\An}$.

\emph{Step 4: conclusion.}  Taking the supremum over all competitors
$\tau$ gives $\sup\{\tau\}\le u$ pointwise, while Step~2 gives the
reverse inequality because $u$ is itself a competitor.  Hence the raw
supremum equals $u$; as $u$ is upper semicontinuous, its regularization is
again $u$, and
\[
                         P_L(\phi_D)=u\in\PSH_X(L).
\tag{40.29}\label{eq:bj-smooth-obstacle-envelope}
\]

\Cref{prop:bj-smooth-obstacle-interface} now gives EP in every real
numerical class.  Nothing in the final comparison discards the common base
locus: there the ideal functions $q_m$, the limit $Q$ and the psh
competitors may take the value $-\infty$ --- the obstacle $\phi_D$ itself
is a PL function, hence bounded and continuous --- and all displayed
inequalities retain their extended-real meaning.  In particular Step~3 is an inequality between two
functions with values in $[-\infty,+\infty)$, and no difference of two
such values is formed; and the identification in Step~4 is an identity of
functions on the whole of $X^{\An}$, not merely on the full-support locus.
\end{proof}

\begin{remark}[Where each hypothesis is used]
\label{rem:ch40-hypothesis-map}
Smoothness of $X$ is used exactly twice, and both times through
$S=X\times\mathbf A^1$ rather than through $X$ itself.  It is used to
make $S$ regular, so that multiplier ideals have the valuative description
\Cref{thm:appC-valuative-multiplier} with $A_S>0$ --- whence
$J_m\subseteq B_m$ --- and so that $S$ is $F$-finite regular, whence
Kunz's criterion and the trace description
\eqref{eq:appC-test-trace}.  And it is used in the subadditivity theorem, in
characteristic zero through the smoothness of the diagonal of $S\times S$
and in characteristic $p$ through the flatness of Frobenius.  It is
\emph{not} used to resolve $\cY$: the model $\cY$ is only normal, and no
statement about $\cY$ beyond $\mu_*\cO_{\cY}=\cO_S$ is needed.

Projectivity of $X$ is used for the compactness of $X^{\An}$, in the
Dini argument of \Cref{prop:ch40-envelope-comparison}, and for the coherence
of $\rho_*$ and the vanishing theorems.  Integrality is used throughout,
in particular for generic finiteness.  The characteristic is used only in
\Cref{sec:ch40-engines} and \Cref{sec:ch40-stages}, and \emph{only} to
choose which of the two engines is installed.  Algebraic closedness of
$k$ is used for the Vandermonde argument of
\Cref{prop:ch40-engine-package}\,\cref{itm:ch40-eng:equivariant}, where an
infinite field would suffice, and through the standing conventions of
\Cref{part:bj}.
\end{remark}

\begin{example}[The envelope of the double point obstacle]
\label{ex:ch40-p2-envelope}
Return to \Cref{ex:ch40-p2-obstacle} and \Cref{ex:ch40-p2-engines}.  All the
ingredients have been computed, so the theorem can be read off:
\[
 \begin{gathered}
 u=\lim_mu_m=-\bigl(1+\min\{s,1\}\bigr)=Q,\\
 P_{\cO(1)}\bigl(\max\{2\log|I_p|,-2\}\bigr)
 =\max\bigl\{\log|I_p|,-1\bigr\}-1 .
 \end{gathered}
\]
Three checks confirm that this is the right answer.

\emph{It is psh.}  The function $\max\{\log|I_p|,-1\}=\phi_{-\Xi}$ is
$c_1(\cO(1))$-psh by the nefness computation in
\Cref{ex:ch40-p2-obstacle}\,(i), and adding the constant $-1$ preserves
this by \Cref{thm:bj-basic-closure}\,(1).

\emph{It is a competitor.}  $-(1+\min\{s,1\})\le-2\min\{s,1\}$, with
equality exactly on $\{s\ge1\}$.  So the envelope touches the obstacle
precisely at the semivaluations with $v(I_p)\ge1$, and is strictly below
it elsewhere, including at the trivial valuation, where the envelope takes
the value $-1$ and the obstacle the value $0$.

\emph{It is maximal.}  If $\tau\in\PSH_X(\cO(1))$ satisfies
$\tau\le\phi_D$, then $\tau\le u$ by Step~3 of the proof of
\Cref{thm:bj-smooth-envelope}.  Independently, no larger candidate of the
same shape works: for $c>0$ the function $\phi_{-\Xi}-1+c$ fails
$\le\phi_D$ at any point with $s\ge1$, where it equals $-2+c>-2$.

Finally, note the structure of the answer.  The obstacle demanded a pole of
order two at $p$; the class $\cO(1)$ can support a pole of order at most
one; the envelope takes the largest admissible pole, namely one, and pays
for the mismatch with the additive constant $-1$.  The constant is not
visible in the divisor $D$ and is produced by the base ideals: it is the
$\varpi^m$ factor in $J_m=\varpi^m\mathfrak m^m$.  This is the second
grading of \Cref{prop:ch40-two-gradings-here} doing its work.
\end{example}

\section{Comparison of the two proofs}
\label{sec:ch40-comparison}

The two proofs have the same envelope skeleton but different ideal
engines.  In characteristic zero, asymptotic multiplier ideals provide the
subadditivity, uniform generation, and valuation approximation used in
\cite{BJ22b}.  In characteristic $p>0$, the corresponding inputs are:
\[
\begin{array}{c|c}
\text{required operation} & \text{positive-characteristic replacement}\\
\hline
\text{subadditive ideal stages}
  & \tau(J_\bullet^{m\ell})\subset\tau(J_\bullet^m)^\ell\\
\text{uniform generation}
  & \text{Frobenius trace plus the fixed twist }A=K_X+dH\\
\text{valuation control}
  & w(\tau(I^t))\ge t\,w(I)-A_S(w).
\end{array}
\]
The class carried by $u_m$ is exactly $L+A/m$, not merely a class
within an unspecified $O(m^{-1})$ neighborhood of $L$.  The error in a
fixed divisorial value is exactly bounded by $A_S(w)/m$.  Keeping these
two errors separate is what makes the limiting argument transparent.

\subsection{The full comparison}
\label{sec:ch40-comparison-table}

The following table records, row by row, every step of the argument and what
implements it in each characteristic.  Rows marked $\ast$ are the four
places where the implementations have no proof step in common.

\begin{center}
\begin{tabular}{p{0.23\textwidth}p{0.34\textwidth}p{0.34\textwidth}}
\toprule
step & $\operatorname{char}k=0$ & $\operatorname{char}k=p>0$\\
\midrule
reduction to a rational obstacle &
\multicolumn{2}{p{0.70\textwidth}}{identical:
\Cref{prop:bj-smooth-obstacle-interface}, using only
\Cref{thm:bj-basic-closure} and
\Cref{thm:general-obstacle-criterion}}\\
\addlinespace
graded sequence $J_\bullet$ &
\multicolumn{2}{p{0.70\textwidth}}{identical:
\Cref{lem:ch40-J-basic}; no resolution of $\cY$ in either case}\\
\addlinespace
ambient scheme &
\multicolumn{2}{p{0.70\textwidth}}{identical: $S=X\times\mathbf A^1$,
smooth over $k$; in characteristic $p$ also $F$-finite}\\
\addlinespace
stage ideal $B_m$$^\ast$ &
$\mathcal J_S(J_\bullet^m)$, defined by a log resolution and a floor
function &
$\tau_S(J_\bullet^m)$, defined by stabilized Frobenius traces\\
\addlinespace
containment $J_m\subseteq B_m$ &
valuative criterion \eqref{eq:appC-valuative-multiplier} with
$A_S>0$ &
trace description \eqref{eq:appC-test-trace}\\
\addlinespace
subadditivity$^\ast$ &
\cite{DEL00}: external product on $S\times S$, restriction to the
diagonal, Nadel vanishing &
\cite{MustataTest}: direct manipulation of Frobenius roots and rounding\\
\addlinespace
fixed twist$^\ast$ &
relative Kawamata--Viehweg vanishing on a log resolution
(\Cref{prop:ch40-multiplier-uniform-generation}) &
Frobenius trace, projection formula, Serre vanishing
(\Cref{prop:bj-smooth-uniform-generation})\\
\addlinespace
role of $K_X$ in $A$ &
$K_T$ is forced by the vanishing theorem &
$\omega_S$ is forced by the trace\\
\addlinespace
why one $A$ suffices$^\ast$ &
the fractional part $\{\tfrac mqF\}$ has coefficients in $[0,1)$ for
every $m$ &
the residual coefficient
\eqref{eq:bj-smooth-residual-coefficient} lies in
$\{-(q-1),\ldots,0\}$ while $p^e(d-i)\to\infty$\\
\addlinespace
stage class &
\multicolumn{2}{p{0.70\textwidth}}{identical: exactly
$c_1(L)+m^{-1}c_1(A)$, nef-decreasing to $c_1(L)$}\\
\addlinespace
valuation error$^\ast$ &
$A_S(v)$ from the relative canonical divisor on a log resolution &
$A_S(w)$ from Canton's Cartier-algebra discrepancy \cite{Canton}\\
\addlinespace
passage to the limit &
\multicolumn{2}{p{0.70\textwidth}}{identical:
\Cref{def:bj-psh}, or clauses (4) and (5) of
\Cref{thm:bj-basic-closure}}\\
\addlinespace
identification &
\multicolumn{2}{p{0.70\textwidth}}{identical:
\Cref{prop:ch40-envelope-comparison} and
\Cref{prop:ch39-determination-package}}\\
\bottomrule
\end{tabular}
\end{center}

\subsection{What each proof cannot do in the other characteristic}
\label{sec:ch40-comparison-obstructions}

\emph{The characteristic-zero proof in characteristic $p$.}  It fails at
two rows.  Log resolutions of $(S,J_q)$ are not known to exist for
$\dim S\ge4$, so \eqref{eq:appC-multiplier-resolution} has no meaning; and
even in dimension three, where resolutions are available, Kawamata--Viehweg
vanishing is false, so Step~4 of
\Cref{prop:ch40-multiplier-uniform-generation} breaks.  This is precisely
why \Cref{thm:bj22-nef-envelope} carries the hypothesis
``$\operatorname{char}k=0$ or $\dim X\le2$'': in the surface case
$\dim S=3$, resolutions exist, and the missing vanishing is supplied by
the positive-characteristic technology of \cite{GJKM19} rather than by
multiplier ideals.

\emph{The characteristic-$p$ proof in characteristic zero.}  It fails at
the first row of the engine.  There is no Frobenius, so
\eqref{eq:appC-test-trace} does not define anything; the trace surjection
\eqref{eq:bj-smooth-trace-quotient} has no analogue; and the mechanism by
which a bounded residual coefficient is traded against an unboundedly
growing ample coefficient has no counterpart, because there is no
$p^e$ to grow.  Test ideals are not a characteristic-free construction
that happens to be usually written in characteristic $p$; they are a
construction \emph{about} $p$.

\emph{The two engines are not the same sheaf.}  \Cref{ex:ch40-cusp} exhibits
a graded sequence and an exponent for which
$\mathcal J$ is the unit ideal and $\tau$ is not.  The comparison
theorem \Cref{thm:appC-hara-yoshida} relates them only after spreading out
and excluding primes, as \Cref{warn:ch40-no-reduction} explains.

\begin{remark}[Multiplier ideals versus test ideals]
\label{rem:ch40-multiplier-versus-test}
The proof does not assert that multiplier ideals and test ideals agree
after reduction modulo $p$, nor does it import a characteristic-zero
resolution theorem into positive characteristic.  It uses only the
properties stated in \eqref{eq:bj-test-ideal-properties}, the trace
surjection \eqref{eq:bj-smooth-trace-quotient}, and Canton's estimate.
These are sufficient for the envelope theorem.
\end{remark}

\begin{remark}[What the parallel buys, beyond uniformity of statement]
\label{rem:ch40-parallel-value}
It is tempting to regard the two-engine presentation as an expository
device: one proof, written twice.  It is more than that, for two reasons
that matter downstream.

First, the parallel isolates the interface.  Everything outside
\Cref{sec:ch40-engines,sec:ch40-stages} is characteristic-free, and the
interface consists of exactly four statements, listed in
\Cref{prop:ch40-engine-package}.  Any future construction supplying those
four statements --- for a scheme that is neither smooth in characteristic
zero nor $F$-finite regular in characteristic $p$, say --- would extend
\Cref{thm:bj-smooth-envelope} verbatim.  Recording the interface is how a
proof becomes reusable.

Second, the parallel is what makes the rest of Part~IV possible.
\Cref{chap:bj-alterations} reduces the normal case to the smooth case in
\emph{both} characteristics simultaneously, using
\Cref{chap:bj-finite-descent} and de Jong's alterations; had the smooth
case been available only in characteristic zero, the alteration argument
would produce nothing in characteristic $p$, and
\Cref{chap:bj-universal-homeomorphisms} would have no smooth input to
descend from.  The chapter is therefore not two proofs of one theorem but
one proof of the theorem that the later chapters need.
\end{remark}

\section{Safeguards, failures, and checkpoints}
\label{sec:ch40-audit}

\subsection{The extended-real audit}
\label{sec:ch40-extended-real}

\begin{proposition}
\label{prop:ch40-safeguards}
In the proofs of \Cref{prop:bj-smooth-obstacle-interface},
\Cref{prop:ch40-Q-limit}, \Cref{prop:ch40-envelope-comparison},
\Cref{prop:ch40-engine-package},
\Cref{prop:ch40-multiplier-uniform-generation},
\Cref{prop:bj-smooth-uniform-generation},
\Cref{prop:bj-smooth-divisorial-identification} and
\Cref{thm:bj-smooth-envelope}, every extended-real operation is one of the
following, and each is legitimate for the reason stated.
\begin{enumerate}
\item\label{itm:ch40-safe:ideal}
 \emph{Ideal values.}  For a nonzero coherent ideal $I\subseteq\cO_S$ and
 $v\in X^{\An}$, the number $\sigma(v)(I)$ lies in $[0,+\infty]$, and
 it is $+\infty$ exactly when the support of $\sigma(v)$ contains $I$,
 by \Cref{prop:ch32-tropical-identities}.  Hence
 $\phi_I$ takes values in $[-\infty,0]$.  All identities used, namely
 $\phi_{II'}=\phi_I+\phi_{I'}$ and
 $\phi_{I+I'}=\max\{\phi_I,\phi_{I'}\}$, are identities in
 $[-\infty,0]$, by \eqref{eq:ch34-log-ideal-semiring}; no difference is
 formed.
\item\label{itm:ch40-safe:identity}
 \emph{The obstacle identity.}  In
 \eqref{eq:bj-smooth-obstacle-identity} all three terms lie in
 $[-\infty,0]$ and the identity is asserted as an identity of sums.  It is
  never solved for $\phi_D$ at a point where
 $\phi_{\mathfrak a_m}=-\infty$; on $X^{\divis}$ all three terms are
 finite, and that is where it is solved.
\item\label{itm:ch40-safe:squeeze}
 \emph{The valuation squeeze.}  In
 \eqref{eq:bj-smooth-canton-squeeze} all three quantities are finite real
 numbers, because $w=\sigma(v)$ is a valuation of $K(S)$ for
 $v\in X^{\divis}$ and $J_m,B_m$ are nonzero.  Division by $m$ and
 negation are therefore ordinary arithmetic.
\item\label{itm:ch40-safe:limits}
 \emph{Monotone limits.}  The limits $q_m\uparrow Q$ and
 $u_m\downarrow u$ are monotone limits in $\Rbar$, defined pointwise
 without hypothesis by ledger entry~(1) of
 \Cref{sec:extended-arithmetic}.  Both are limits along the divisibility
 order, which is directed but not the usual order on $\mathbf N$; the
 distinction is recorded in
 \Cref{prop:ch40-Q-limit}\,\cref{itm:ch40-Q:increasing}.
\item\label{itm:ch40-safe:dini}
 \emph{The Dini step.}  In Step~2 of
 \Cref{prop:ch40-envelope-comparison} the sets
 $\{\tau_i<h+\varepsilon\}$ are open because $\tau_i$ is usc with values
 in $[-\infty,+\infty)$ and $h$ is real valued and continuous; the
 inequality $\tau_i\le h+\varepsilon$ is an inequality in
 $[-\infty,+\infty)$.  No difference $\tau_i-h$ is formed at a point
 where $\tau_i=-\infty$.
\item\label{itm:ch40-safe:competitor}
 \emph{The competitor $\chi_{i,\delta}$.}  The expression
 $(1-\delta)\tau_i+\delta(\log|s_E|+C)-\varepsilon$ is a sum of two
 functions with values in $[-\infty,+\infty)$ and a real constant, hence
 defined everywhere by ledger entry~(2); it is evaluated at a divisorial
 point, where both summands are finite, only in Step~5.
\item\label{itm:ch40-safe:comparison}
 \emph{The two comparisons.}  Both applications of
 \Cref{prop:ch39-determination-package}\,\cref{itm:ch39-det:compare} in the
 proof of \Cref{thm:bj-smooth-envelope} have the psh function as the
 \emph{smaller} one: $u\le\phi_D$ in Step~2, $\tau\le u$ in Step~3.  The
 asymmetry is essential and is recorded in
 \Cref{warn:ch39-determination-hypotheses}.
\item\label{itm:ch40-safe:sup}
 \emph{The final supremum.}  In Step~4 of
 \Cref{thm:bj-smooth-envelope} the supremum over competitors is bounded
 above pointwise by the fixed function $u$, so it is not a supremum of an
 unbounded family; and it is shown to \emph{equal} $u$, which is usc, so
 that the regularization changes nothing.  At no point is a supremum of psh
 functions asserted to be psh before this identification.
\end{enumerate}
\end{proposition}

\begin{proof}
Each item was verified where it occurs; the list is a record, and the
justifications are those given in the corresponding proofs, together with
\Cref{lem:extended-real-safe} and the ledger of
\Cref{sec:extended-arithmetic}.
\end{proof}

\subsection{What fails if a safeguard is dropped}
\label{sec:ch40-safeguard-failures}

\emph{Drop the Veronese normalization.}  If $J_m=0$ for some $m$, then
$\phi_{J_m}\equiv-\infty$, the ideal $B_m$ is not defined by
\Cref{def:ch40-stage-ideals}, and \eqref{eq:bj-smooth-gf-lower-bound} has no
content.  The normalization
\Cref{lem:ch40-J-basic}\,\cref{itm:ch40-J:nonzero} is what excludes this.
The restriction of the index set is mild, and it is a restriction at small
$m$ only: for $L$ big, Kodaira's lemma writes $m_1L\sim A_0+E$ with
$A_0$ ample and $E\ge0$ for some $m_1$, whence
$H^0(X,mL)\neq0$ for all $m\gg0$, so only finitely many indices can be
excluded.  The point of the normalization is therefore to fix one $m_1$
with $H^0(X,m_1L)\ne0$ and run the divisibility net through its
multiples, not to avoid infinitely many vanishing stages.

\emph{Drop the fixed twist.}  Suppose one allowed $A$ to depend on $m$,
say $A_m=K_X+d_mH$ with $d_m\to\infty$, which is what a naive
application of Serre vanishing produces.  Then the stage classes are
$c_1(L)+m^{-1}c_1(A_m)$, and $m^{-1}c_1(A_m)$ need not tend to $0$;
if $d_m\sim m$ it tends to $c_1(H)$, and the limit $u$ would at best
be $(c_1(L)+c_1(H))$-psh.  The conclusion of
\Cref{thm:bj-smooth-envelope} would then be false as stated, since the
envelope property is a statement about a specific class.  The whole point of
\Cref{prop:ch40-multiplier-uniform-generation} and
\Cref{prop:bj-smooth-uniform-generation} is that the Serre threshold is
allowed to depend on $m$ while the \emph{twist} is not.

\emph{Drop subadditivity.}  Without
\Cref{prop:ch40-engine-package}\,\cref{itm:ch40-eng:subadd} the family
$(u_m)$ need not be monotone, and there is no reason for a limit to exist,
let alone to be psh.  One might hope to recover monotonicity by passing to
$\inf_{m'\ge m}u_{m'}$; but an infimum of Fubini--Study functions is not a
Fubini--Study function, and the resulting net would not satisfy
\Cref{def:bj-psh}(3).  The direction of the inclusion matters too:
$B_m^{\,\ell}\subseteq B_{m\ell}$, which is the inclusion one might expect
from gradedness, would give $u_{m\ell}\ge u_m$ and an increasing net,
returning the problem to its starting point.

\emph{Drop the valuation inequality.}  Then one still obtains a psh function
$u$ with $Q\le u\le0$, but no reason to believe $u\le\phi_D$; indeed
$u=0$ satisfies every other requirement when $c_1(L)$ is nef, and is not
the envelope.  \Cref{ex:ch40-p2-envelope} is a case in point: there
$u=Q=\phi_{-\Xi}-1$, while the function $0$ is also
$c_1(\cO(1))$-psh and exceeds the obstacle at every point with $s>0$.

\emph{Drop the restriction to divisorial points in Step~3.}  The inequality
$\tau\le Q$ is proved on $X^{\divis}$ and is false in general on
$X^{\An}$ before the identification is complete: at a point $x$ of
proper support the function $Q$ may be $-\infty$ while $\tau$ is
finite, since $Q$ is a supremum of Fubini--Study functions all of which may
vanish on the same subvariety.  \Cref{ex:bj-common-base-locus} exhibits the
configuration.  What repairs the argument is not a better estimate at $x$
but \Cref{prop:ch39-determination-package}, which transfers the inequality
from $X^{\divis}$ using plurisubharmonicity.

\emph{Drop smoothness of $X$.}  Then $S$ is not regular.  Multiplier
ideals and test ideals do exist in broader singular settings, after one
introduces boundaries, $\mathbf Q$-Gorenstein hypotheses or Cartier
algebras as appropriate; what is unavailable is the boundary-free
four-property package used here.  In characteristic zero the strict
discrepancy input behind $J_m\subseteq B_m$ is lost, while in positive
characteristic the proof can no longer invoke Frobenius flatness, the
rank-one trace presentation and the resulting uniform generation theorem.
The singular case is not treated by silently weakening these steps; it is
treated by
\Cref{chap:bj-finite-descent,chap:bj-alterations,chap:bj-universal-homeomorphisms,chap:bj-subintegral-normalization},
which transport the conclusion from a smooth alteration.  This is why the
smooth case has to be settled first and in all characteristics.

\subsection{Proof checkpoints}
\label{sec:ch40-checkpoints}

\begin{enumerate}
\item The base-ideal identity is
      $J_m\cO_{\mathcal Y}=\mathfrak a_m\cO_{\mathcal Y}(mD)$;
      the sign in \eqref{eq:bj-smooth-obstacle-identity} follows only
      after evaluating ideals and negating.
\item The residual coefficient in
      \eqref{eq:bj-smooth-residual-coefficient} lies in a finite interval
      depending on $q$, while the positive coefficient
      $p^e(d-i)$ tends to infinity.  This is the exact Serre-vanishing
      ledger.
\item The twist $A$ is fixed before $m$ varies.  Otherwise
      $c_1(L+A/m)\to c_1(L)$ would not follow.
\item The lower bound by $q_{m_0}$ proves generic finiteness of the
      decreasing limit.
\item Canton's discrepancy is finite for each Gauss extension used here;
      no uniform bound over all divisorial points is required.
\item The final comparison uses full-support divisorial points to avoid
      the augmented base locus, and retains $-\infty$ on the common
      base locus.
\item The two gradings are kept apart: the degree $m$ is the exponent of
      $\cL$ and the denominator of the normalization, while the
      $\mathbf G_m$-weight $k$ is the exponent of $\varpi$ and enters
      as the Fubini--Study weight $-k/m$.
\item The stage classes are exactly $c_1(L)+m^{-1}c_1(A)$, they decrease
      in the nef order, and the limit is taken by a clause of
      \Cref{thm:bj-basic-closure} that permits the class to move.
\item Every appeal to the envelope property in this chapter is either the
      conclusion being proved or an explicit hypothesis about a
      \emph{different}, already treated class, as in
      \Cref{lem:ch40-ep-nef-limit}.
\end{enumerate}

\section{Chapter dictionary}
\label{sec:ch40-dictionary}

The correspondences established or used in this chapter are the following.
Each row records a statement together with the hypotheses under which it was
proved.  Rows one and two are the translation of an obstacle into ideals;
rows three and four are the engine; rows five and six are the two
$O(m^{-1})$ errors; row seven is the conclusion.

\begin{center}
\begin{tabular}{p{.28\textwidth}p{.32\textwidth}p{.29\textwidth}}
\toprule
analytic datum & algebraic datum & hypothesis\\
\midrule
a rational PL obstacle $\phi_D\le0$ &
a vertical divisor $D\le0$ on a model $\mu:\cY\to X\times\mathbf A^1$ &
$X$ integral projective; \Cref{thm:tc-pl-correspondence} and Step~4 of
\Cref{prop:bj-smooth-obstacle-interface}\\
the Fubini--Study envelope $Q_L(\phi_D)$ &
$\lim_mm^{-1}\phi_{J_m}$ for the graded sequence $J_\bullet$ of
\eqref{eq:bj-smooth-ambient-ideal} &
$L$ big; \Cref{prop:ch40-Q-limit}\\
a decreasing stage $u_m$ &
$m^{-1}\phi_{B_m}$ with
$B_{m\ell}\subseteq B_m^{\,\ell}$ &
$S$ smooth over $k$; $F$-finite if
$\operatorname{char}k=p$
(\Cref{prop:ch40-engine-package})\\
the stage as a Fubini--Study function &
relative generation of $B_m(mL+A)$, $A=K_X+dH$, $d\ge n+1$ &
\Cref{prop:ch40-multiplier-uniform-generation} in characteristic zero,
\Cref{prop:bj-smooth-uniform-generation} in characteristic $p$\\
class error &
$\theta_m-c_1(L)=m^{-1}c_1(A)$, nef-decreasing to $0$ &
$A$ independent of $m$
(\Cref{lem:ch40-drift-nef-monotone}, \Cref{ex:ch40-drift-real})\\
value error at $v\in X^{\divis}$ &
$0\le u_m(v)-Q(v)\le A_S(\sigma(v))/m$ &
$\sigma(v)$ divisorial with $A_S(\sigma(v))=A_X(v)+1$
(\Cref{lem:ch40-gauss-log-discrepancy},
\Cref{prop:bj-smooth-divisorial-identification})\\
$P_L(\phi_D)\in\PSH_X(c_1(L))$, hence $\operatorname{EP}_X(\theta)$ &
$u=\lim_mm^{-1}\phi_{B_m}$ &
$X$ smooth projective, any characteristic, any real class
(\Cref{thm:bj-smooth-envelope})\\
\bottomrule
\end{tabular}
\end{center}

Five separations should be carried forward.

\emph{Increasing versus decreasing.}  The base-ideal net $(q_m)$ increases
and the stage net $(u_m)$ decreases; they have the same limit on
$X^{\divis}$ and are compared through
\Cref{prop:ch40-engine-package}\,\cref{itm:ch40-eng:sandwich}.  Only the
decreasing net is admissible in \Cref{def:bj-psh}.  The whole function of
the ideal engine is to convert one into the other, and it does so by an
order-reversing operation on ideals, not by any manipulation of functions.

\emph{Degree versus $\mathbf G_m$-weight.}  In $J_m=\bigoplus_k\mathfrak
b_{m,k}\varpi^k$, the index $m$ is a degree and the index $k$ a weight;
they scale differently, they produce different kinds of error, and only the
first appears as a denominator.  \Cref{prop:ch40-two-gradings-here} records
the rules, \Cref{ex:ch40-p2-obstacle}\,(iv) exhibits $m+1$ distinct weight
pieces in degree $m$, and \Cref{warn:ch40-two-gradings} records what goes
wrong if the weight is suppressed.

\emph{Class error versus value error.}  Both are $O(m^{-1})$; the first
lives in $N^1(X)_{\mathbf Q}$ and is removed by a limit clause of
\Cref{thm:bj-basic-closure}, the second is a real number attached to a
divisorial point and is removed by dividing by $m$.
\Cref{ex:ch40-drift-real} shows that each can be nonzero when the other
vanishes.

\emph{Multiplier ideal versus test ideal.}  The two engines satisfy the same
four properties and are different sheaves; \Cref{ex:ch40-cusp} separates
them, and \Cref{warn:ch40-no-reduction} explains why the comparison theorem
\Cref{thm:appC-hara-yoshida} is not a bridge.  The corresponding separation
on the side of proofs is \Cref{rem:ch40-subadditivity-two-proofs}.

\emph{Envelope property as conclusion versus as hypothesis.}  The only
occurrence of $\operatorname{EP}$ as a hypothesis in this chapter is in
\Cref{lem:ch40-ep-nef-limit}, where it is assumed for the approximating
classes $\theta_j$ and discharged in Step~2 of
\Cref{prop:bj-smooth-obstacle-interface}.  Everywhere else it is the
conclusion.  Reading \Cref{thm:bj-final-envelope} back into
\Cref{prop:ch40-envelope-comparison} or into Step~3 of
\Cref{thm:bj-smooth-envelope} would make
\Crefrange{chap:bj-cartier-realization}{chap:bj-envelope-consequences}
circular (\Cref{warn:ch40-firewall}).

\section{Source note}
\label{sec:ch40-source-note}

The proof follows the smooth-case argument in the supplied envelope
manuscript, which is the source here for the arbitrary-dimensional
positive-characteristic theorem.  The published addendum \cite{BJ22b}
supplies the characteristic-zero obstacle reduction but proves its smooth
theorem only under the alternative
$\operatorname{char}k=0$ or $\dim X\le2$.  Asymptotic test ideals,
subadditivity and the fixed-twist global-generation template are from
\cite[Proposition~3.1 and Theorem~4.1]{MustataTest}; the discrepancy estimate
is from \cite{Canton}.  The relative fixed-twist computation and the
proper-support bookkeeping have been written out because both are necessary
for the stated characteristic-uniform conclusion.

The precise external attributions are as follows.  The historical smooth
cases are \cite[Theorem~5.20]{BJ22}, for nef classes with
$\operatorname{char}k=0$ or $\dim X\le2$, and \cite[Theorem~A]{BJ22b},
which removes the nefness restriction under the same
characteristic-or-dimension alternative; both are recorded in
\Cref{chap:bj-envelope-history} as \Cref{thm:bj22-nef-envelope} and
\Cref{thm:bj-addendum-all-classes}, and neither is used as an input here.
The positive-characteristic surface case, which is where the alternative
``$\dim X\le2$'' in those statements comes from, is
\cite[\S4 and Theorem~4.6]{GJKM19}.  Multiplier ideals, local vanishing and
asymptotic multiplier ideals are \cite[Chapters~9 and~11]{PAG2}, imported
through \Cref{sec:appC-algebraic-multiplier} and
\Cref{sec:appC-asymptotic-multiplier}; subadditivity is
\cite[Theorem~2.3]{DEL00}; the valuative formalism for graded sequences is
\cite{JM12}.  Kawamata--Viehweg vanishing, used in its relative form over
$\mathbf A^1$ in \Cref{prop:ch40-multiplier-uniform-generation}, is
\cite[\S4.3]{PAG1}, and local vanishing is \cite[Theorem~9.4.1]{PAG2};
Castelnuovo--Mumford regularity, used relative to a noetherian base ring, is
\cite[\S1.8]{PAG1}.  Resolution of singularities in characteristic zero
enters only through \Cref{thm:appB-resolution}, which is
\cite{Hironaka1964I}.  Generalized and asymptotic test ideals, their
conventions and their Brian\c con--Skoda property are
\cite{HY03} and \cite{ST12}; right continuity of $c\mapsto\tau(f^c)$,
used in \Cref{ex:ch40-cusp}, is \cite[Theorem~6.6]{ST12}; the comparison
across characteristic, quoted only to be excluded, is
\cite[Theorem~6.8]{HY03}.  The description of the augmented base locus
mentioned in \Cref{rem:ch40-off-base-locus} is \cite{ELMNP05}, and the
identity \eqref{eq:bj-smooth-off-base} itself is
\cite[Lemma~1.5]{BJ22b}; neither is used in the proof.

The $F$-pure threshold of the cusp used in \Cref{ex:ch40-cusp} --- namely
$\operatorname{fpt}(x^2+y^3)=5/6$ for $p\equiv1\bmod3$ and
$5/6-1/(6p)$ for $p\equiv2\bmod3$, $p\ge5$ --- is computed in
\cite[Example~4.3]{MTW05}.  The example uses only the single value
$\operatorname{fpt}_5(x^2+y^3)=4/5$; right continuity is separately
attributed above to \cite[Theorem~6.6]{ST12}.

Seven things in this chapter are not in the sources in this form and are
supplied here.  \Cref{lem:ch40-ep-nef-limit} and
\Cref{lem:ch40-rational-approximation} make the reduction from real classes
to big rational ones explicit, including an elementary proof that a
pseudoeffective class plus an ample class is big; the corresponding step in
\cite{BJ22b} is a reference to \cite[Lemma~5.9]{BJ22}.
\Cref{lem:ch40-equivariant-ideal-fs} is the version of
\Cref{prop:ch34-fs-flag-ideal} for equivariant ideals that are not trivial
over $\mathbf G_m$, which is the case actually needed because $L$ is
only big; the flag-ideal statement of \Cref{chap:bj-test-configurations}
does not apply verbatim.  \Cref{prop:ch40-Q-limit} proves, rather than
cites, that the base-ideal stages exhaust the Fubini--Study envelope, the
membership test being run through the component valuations of
\Cref{lem:ch33-gauss-component}.  \Cref{prop:ch40-envelope-comparison} replaces
the appeal to \eqref{eq:bj-smooth-off-base} by a self-contained argument
valid at every divisorial point, using a Kodaira decomposition and a Dini
step; this removes the augmented base locus from the proof entirely.
\Cref{prop:ch40-multiplier-uniform-generation} writes out the
characteristic-zero relative generation statement in the same relative form
as its positive-characteristic counterpart, so that the two can be compared
line by line; the source states the absolute version.
\Cref{lem:ch40-gauss-log-discrepancy} computes
$A_S(\sigma(v))=A_X(v)+1$, which is what makes
\Cref{rem:ch40-no-uniform-discrepancy} a checkable statement about the
failure of uniformity.  Finally the worked computations
\Cref{ex:ch40-p2-obstacle}, \Cref{ex:ch40-p2-engines},
\Cref{ex:ch40-cusp}, \Cref{ex:ch40-monomial-valuation},
\Cref{ex:ch40-drift-real} and \Cref{ex:ch40-p2-envelope} are new.  They are
included because the two claims this chapter exists to establish --- that
the same architecture runs in every characteristic, and that the two
$O(m^{-1})$ errors are separately controlled --- are exactly the claims
that a plausible argument can appear to establish without doing so.  A
reader who has computed $J_m=\varpi^m\mathfrak m^m$ and
$B_m=\varpi^m\mathfrak m^{m-2}$ by hand, and has seen the two engines
separate on the cusp, will not later mistake the parallel for an identity.
